\subsection{Characterizations of the interval}

Lemma 6. --- Let D be a countable totally ordered set, not reduced to a single element, having a least and a greatest element. Suppose that D is without gaps (E, III, p.~73, exerc. 19), that is, every open interval {]}x, y{[}, where x and y are elements of D such that x \textless{} y, is nonempty. There then exists an isomorphism of ordered sets from I ∩ Q onto D.

Let \(a\) be the smallest element and \(b\) the largest element of D. By hypothesis, we have \(b \neq a\) and \(]a\) , \(x[\neq \emptyset\) for all \(x \in \mathrm{D} - \{a\}\) . The set \(\mathrm{D} - \{a\}\) is totally ordered, is nonempty, and has no least element; it is therefore infinite (E, III, p.~34, cor. 1 of prop. 3). The sets \(\mathbf{I} \cap \mathbf{Q}\) and D, infinite and countable, are equipotent to \(\mathbf{N}\) .

Let us choose bijections \(n \mapsto a_{n}\) and \(n \mapsto b_{n}\) from N onto \(I \cap Q\) and D respectively, such that \(a_{0} = 0\) , \(a_{1} = 1\) , \(b_{0} = a\) , \(b_{1} = b\). There exists a unique strictly increasing map \(f: I \cap Q \to D\) possessing the following properties: one has \(f(0) = a\) and \(f(1) = b\) ; for \(n \geqslant 2\) , we have \(f(a_{n}) = b_{m}\) , where m is the smallest natural integer for which the map from \(\{a_{0}, \ldots, a_{n}\}\) into D that coincides with f in \(\{a_{0}, \ldots, a_{n-1}\}\) and maps \(a_{n}\) to \(b_{m}\) is strictly increasing. These properties indeed define \(f(a_{n})\) by induction on n, the existence of the integer m being ensured by the fact that D has no gaps.

Since the map \(f\) is strictly increasing and \(\mathbf{I} \cap \mathbf{Q}\) is totally ordered, \(f\) defines an isomorphism of ordered sets from \(\mathbf{I} \cap \mathbf{Q}\) onto its image (E, III, p.~14, prop. 11). It remains for us to prove that \(f\) is surjective. To do this, let us prove by induction that \(b_{m}\) belongs to the image of \(f\) for all \(m \in \mathbf{N}\) .

We have \(b_{0}=f(0)\) and \(b_{1}=f(1)\) . Suppose that \(m\geqslant2\) and that, for \(0\leqslant k\leqslant m-1\) , there exists \(c_{k}\in I\cap Q\) such that \(f(c_{k})=b_{k}\) . We have \(c_{0}=0\) and \(c_{1}=1\) , since f is injective. Consider the smallest integer \(n\in N\) for which \(a_{n}\) does not belong to \(\{c_{0},\ldots,c_{m-1}\}\) and the map g from \(\{c_{0},\ldots,c_{m-1},a_{n}\}\) into D that coincides with f in \(\{c_{0},\ldots,c_{m-1}\}\) and maps \(a_{n}\) to \(b_{m}\) is strictly increasing; such an integer exists because \(I\cap Q\) is without gaps. Let \(f'\) be the map from \(\{a_{0},\ldots,a_{n}\}\) into D that coincides with f in \(\{a_{0},\ldots,a_{n-1}\}\) and maps \(a_{n}\) to \(b_{m}\). Let us prove that it is strictly increasing. Let \(j\in\{0,\ldots,n-1\}\) ; by definition of the integer n, there exists \(k\in\{0,\ldots,m-1\}\) such that \(a_{j}=c_{k}\) , or \(a_{j}<c_{k}\) and \(b_{m}\geqslant f(c_{k})\) , or \(a_{j}>c_{k}\) and \(b_{m}\leqslant f(c_{k})\) . Suppose \(b_{k}<b_{m}\) ; we then have \(g(c_{k})=f(c_{k})=b_{k}<b_{m}=g(a_{n})\) , whence \(c_{k}<a_{n}\) since g is strictly increasing, and then \(a_{j}\leqslant c_{k}<a_{n}\) ; moreover, \(f'(a_{j})=f(a_{j})\leqslant f(c_{k})=b_{k}<b_{m}=f'(a_{n})\) . Similarly, if \(b_{k}>b_{m}\) , it follows that \(a_{n}<c_{k}\leqslant a_{j}\) and \(f'(a_{j})=f(a_{j})\geqslant f(c_{k})=b_{k}>b_{m}=f'(a_{n})\) . Since f is itself strictly increasing, it follows that the map \(f'\) is strictly increasing.

If \(m'\) is the integer such that \(f(a_n) = b_{m'}\) , we have \(m' \leqslant m\) , by definition of f. If one had \(m' < m\) , one would have \(f(a_n) = b_{m'} = f(c_{m'})\) , whence \(a_n = c_{m'}\) , since f is injective, which contradicts the definition of \(a_n\) . Thus, \(m' = m\) and \(f(a_n) = b_m\) , which proves that \(b_m\) belongs to the image of f and completes by induction the proof of the surjectivity of f.

PROPOSITION 16. --- Let E be a totally ordered set not reduced to a single element. Suppose that every subset of E has a supremum and that there exists a countable subset of E that meets every open interval of the form {]}x, y{[}, where x and y are elements of E such that x \textless{} y. There then exists an isomorphism of ordered sets from I onto E.

Since ∅ and E each have a supremum in E, E has a least element a and a greatest element b. These are distinct since E is not reduced to a single element. Let D' be a countable subset of E that meets every open interval of E of the form {]}x, y{[} with x \textless{} y. Set D = D' ∪ \{a, b\}. The set D is totally ordered and has no gaps. By Lemma 6, there exists an order isomorphism f from I ∩ Q onto D.

For every \(t \in I\) , let \(g(t)\) be the upper bound of \(f([0, t] \cap \mathbf{Q})\) in E. For every \(x \in E\) , let \(h(x)\) be the upper bound of \(f^{-1}([a, x] \cap D)\) in I. The maps \(g: I \to E\) and \(h: E \to I\) thus defined are increasing, g coincides with f in \(I \cap Q\) and h coincides with \(f^{-1}\) in D.

We therefore have \(g(h(y)) = y\) for all \(y \in D\) . Let \(x \in E\) . If one had \(g(h(x)) > x\) , the interval {]}x, \(g(h(x))[\) would contain a point y of D and the relations \(g(h(y)) = y < g(h(x))\) would contradict the fact that \(g \circ h\) is increasing. Likewise, one does not have \(g(h(x)) < x\) . We therefore have \(g(h(x)) = x\) , which proves that \(g \circ h\) is the identity map of E. One shows similarly that \(h \circ g\) is the identity map of I. Thus, \(g: I \to E\) and \(h: E \to I\) are mutually inverse isomorphisms of ordered sets.

Remark. --- Let E be a totally ordered set. The set of open intervals of E (bounded or not) is stable under finite intersection. It is a base of a topology \(\mathcal{T}_{0}(\mathrm{E})\) on E (TG, I, p.~91, exerc. 5). The topology \(\mathcal{T}_{0}(\mathbf{I})\) is identical to the topology induced on \(\mathbf{I}\) by that of \(\mathbf{R}\) . It follows that, in prop. 16, every isomorphism of ordered sets from \(\mathbf{I}\) onto E is a homeomorphism of \(\mathbf{I}\) onto the topological space obtained by endowing E with the topology \(\mathcal{T}_{0}(\mathrm{E})\) .

COROLLARY. --- Let R be an equivalence relation in I. The following conditions are equivalent:

\begin{enumerate}
\def\labelenumi{(\roman{enumi})}
\item
  Every equivalence class modulo R is a closed interval of I, distinct from I;
\item
  There exists an increasing surjective map \(u\colon\mathbf{I}\to\mathbf{I}\) such that R is the equivalence relation associated with \(u\) .
\end{enumerate}

Such a map u, when it exists, is continuous and induces, by passage to the quotient, a homeomorphism of I/R onto I.

We shall denote by \(p: I \rightarrow I/R\) the canonical surjection.

Suppose condition (i) is satisfied. For A and B equivalence classes modulo R, write A \textless{} B if one has a \textless{} b for all \(a \in A\) and every \(b \in B\) . In I/R, the relation “A = B or A \textless{} B” is an order relation. Indeed, it is reflexive; it is antisymmetric because one cannot have simultaneously A \textless{} B and B \textless{} A; it is transitive because the relations A \textless{} B and B \textless{} C imply A \textless{} C. If A and B are distinct elements of I/R, they are disjoint closed intervals of I, and one then has either A \textless{} B or B \textless{} A. Endowed with the order relation thus defined, I/R is therefore totally ordered. The map \(p: I \to I/R\) is increasing.

The set I/R is not reduced to a single element, by virtue of (i).

Let F be a subset of I/R. Let us prove that F has a least upper bound in I/R. Put \(F' = \overline{p}^{-1}(F)\) ; denote by a the least upper bound of \(F'\) in I and by A the equivalence class of a under R. Since a majorizes \(F'\) in I, A majorizes F in I/R. Conversely, let \(B \in I/R\) be an upper bound of F; put \(b = \sup(B)\) . Every element of \(F'\) is then majorized by b. We therefore have \(a \leqslant b\) . Since the equivalence classes under R are closed intervals, a belongs to A and b to B. Consequently, we have \(A = p(a) \leqslant p(b) = B\) . This proves that A is the least upper bound of F.

Let A and B be elements of I/R such that A \textless{} B. Let a be the least upper bound of A and b the greatest lower bound of B. Since \(a \in A\) and \(b \in B\) , we have a \textless{} b. The equivalence class under R of any element of {]}a, b{[} is an element of the interval {]}A, B{[} of I/R. Since \(I \cap Q\) meets {]}a, b{[}, its image under p meets {]}A, B{[}.

We have thus proved that the totally ordered set I/R satisfies the hypotheses of prop. 16. There therefore exists an isomorphism of ordered sets \(f \colon \mathbf{I} / \mathrm{R} \to \mathbf{I}\) . The map \(u = f \circ p\) is a surjective and increasing map from I onto I and the equivalence relation associated with u is the relation R; this proves that condition (ii) is satisfied.

Suppose conversely that condition (ii) is satisfied. Let \(u: I \rightarrow I\) be an increasing and surjective map such that R is the equivalence relation associated with u.

Let \(a\) be a point of \(\mathbf{I}\) . The set \(\mathrm{A} = \overline{u}^{-1}(a)\) is an interval of \(\mathbf{I}\) , because the map \(u\) is increasing. Since \(u\) is surjective, A is neither empty nor equal to \(\mathbf{I}\) . Let \(b\) be the least upper bound of A in \(\mathbf{I}\) . We have \(u(b) \geqslant a\) . If \(u(b) > a\) , there exists \(c \in ]a, u(b)[\) and \(d \in \mathbf{I}\) such that \(u(d) = c\) , since \(u\) is surjective. Since \(u\) is increasing and \(a < u(d) < u(b)\) , \(d\) majorizes every element of A and we have \(d < b\) , which contradicts the hypothesis that \(b\) is the supremum of A. We therefore have \(u(b) = a\) , that is, \(b \in \mathrm{A}\) . Similarly, one shows that A contains its lower bound. The set A is therefore a closed interval of \(\mathbf{I}\) distinct from \(\mathbf{I}\) . This proves that condition (i) is satisfied.

We now prove that the map \(u\) is continuous. It suffices for this to show that, for every \(a \in \mathbf{I}\) , the sets \(\overline{u}^{-1}([a, \to[)\) and \(\overline{u}^{-1}([\leftarrow, a[)\) are open. Let \(b\) be the upper bound of \(\overline{u}^{-1}(a)\) ; we have \(u(b) = a\) . If \(x \in \mathbf{I}\) satisfies \(u(x) > a\) , one necessarily has \(x > b\) ; conversely, if \(x > b\) , we have \(u(x) \geqslant a\) and \(u(x) \neq a\) since \(b\) is the least upper bound of \(\overline{u}^{-1}(a)\) . We therefore have \(\overline{u}^{-1}([a, \to[) = ]b, \to[\) , which shows that the inverse image under \(u\) of the interval \(]a, \to[\) is open. One shows similarly that the inverse image of \(]\leftarrow, a[\) is open. It follows that the map \(u\) is continuous.

The map \(v \colon \mathbf{I} / \mathrm{R} \to \mathbf{I}\) deduced from \(u\) passing to the quotient is then continuous and bijective. Since \(\mathbf{I}\) is compact, \(\mathbf{I} / \mathrm{R}\) is quasi-compact (TG, I, p.~62, th. 2) and \(v\) is a homeomorphism (TG, I, p.~63, cor. 2).

PROPOSITION 17. --- Let X be a connected and compact topological space. Let a be a point of X, let U be a nonempty open and closed subset of \(X - \{a\}\).

\begin{enumerate}
\def\labelenumi{\alph{enumi})}
\item
  The closure \(\overline{U}\) of U in X is equal to \(U \cup \{a\}\) and is connected.
\item
  Let \(a'\) be a point of X distinct from \(a\) , let \(U'\) be a nonempty open and closed subset of \(X - \{a'\}\) . If \(a \notin U'\) and \(\overline{U} \cap \overline{U'} \neq \varnothing\) , we have \(\overline{U'} \subset U\) . Conversely, if \(a \in U'\) and \(X \neq U \cup U'\) , we have \(\overline{U} \subset U'\) .
\item
  There exists a point \(b\) of U such that \(X - \{b\}\) is connected.
\end{enumerate}

Let us prove assertion \(a\) ). Let V denote the complement of U in \(X - \{a\}\). Since U is closed in \(X - \{a\}\), V is open in \(X - \{a\}\) and a fortiori in X. We therefore have \(U \subset \overline{\mathrm{U}} \subset X - V = U \cup \{a\}\). Likewise, U is an open subset of X. Since X is connected, U is not closed in X, whence \(\overline{\mathrm{U}} = U \cup \{a\}\). Similarly, we have \(\overline{\mathrm{V}} = V \cup \{a\}\).

Let F and G be disjoint closed subsets of \(\overline{U}\) such that \(\overline{U} = F \cup G\) . Suppose that \(a \in G\) and let us prove that F is empty; we would reason similarly if \(a \in F\) . The set \(G \cup V = G \cup \overline{V}\) is a closed subset of X, disjoint from F, and we have \(X = \overline{U} \cup V = F \cup (G \cup V)\) . Since X is connected and G is nonempty, F is empty. This proves that \(\overline{U}\) is connected.

Let us prove \(b\) ). Suppose that \(a \notin \mathrm{U}'\) and that \(\overline{\mathrm{U}} \cap \overline{\mathrm{U}'} \neq \emptyset\) . By assertion \(a\) ), we have \(\overline{\mathrm{U}} = \mathrm{U} \cup \{a\}\) and \(\overline{\mathrm{U}'} = \mathrm{U}' \cup \{a'\}\) . Since \(a \notin \mathrm{U}'\) and \(a \neq a'\) , the sets U and \(\overline{\mathrm{U}'}\) have a common point. Still according to \(a\) ), \(\overline{\mathrm{U}'}\) is a connected subset of X; since it is contained in \(\mathrm{X} - \{a\}\) and meets the open and closed subset U, we have \(\overline{\mathrm{U}'} \subset \mathrm{U}\) , which is what had to be proved. The second assertion follows from the first by considering the complements in X of \(\overline{\mathrm{U}}\) and \(\overline{\mathrm{U}'}\) respectively.

Let us finally prove assertion \(c\) ). Suppose by contradiction that, for every \(x \in \mathrm{U}\) , the set \(\mathrm{X} - \{x\}\) is not connected, and choose subsets \(\mathrm{U}_x\) and \(\mathrm{V}_x\) , open and closed in \(\mathrm{X} - \{x\}\) , disjoint and nonempty, such that \(\mathrm{X} - \{x\} = \mathrm{U}_x \cup \mathrm{V}_x\) and \(a \in \mathrm{V}_x\) . By assertion \(b\) ), applied to the open and closed subsets \(\mathrm{U}\) , \(\mathrm{U}_x\) of \(\mathrm{X} - \{a\}\) and \(\mathrm{X} - \{x\}\) respectively, we have \(\overline{\mathrm{U}_x} \subset \mathrm{U}\) for all \(x \in \mathrm{U}\) . Let \(x\) and \(y\) be points of U such that \(x \in \mathrm{U}_y\) ; one therefore has \(x \neq y\) . Still by assertion \(b\) ), applied to the open and closed subsets \(\mathrm{U}_x\) and \(\mathrm{U}_y\) of \(\mathrm{X} - \{x\}\) and \(\mathrm{X} - \{y\}\) , the relation \(x \in \mathrm{U}_y\) implies the relation \(\overline{\mathrm{U}_x} \subset \overline{\mathrm{U}_y}\) . Consequently, the relations \(x \in \mathrm{U}_y\) and \(\overline{\mathrm{U}_x} \subset \overline{\mathrm{U}_y}\) are equivalent.

The set of subsets S of U such that \(\bigcap_{x\in\mathrm{S}}\overline{\mathrm{U}_{x}}\neq\varnothing\) is of finite character (E, III, p.~34), since the space X is compact. By E, III, p.~35, th. 1, there exists a maximal subset S of U such that \(\mathrm{C}=\bigcap_{x\in\mathrm{S}}\overline{\mathrm{U}_{x}}\) is not empty. Let \(c\) be a point of C. For every element \(x\) of S, we have \(c\in\overline{\mathrm{U}_{x}}\) , whence \(c\in\mathrm{U}\) ; then \(\overline{\mathrm{U}_{c}}\subset\overline{\mathrm{U}_{x}}\) . Consequently, one has \(\overline{\mathrm{U}_{c}}\subset\mathrm{C}\) ; then, by maximality of S, \(\mathrm{C}=\overline{\mathrm{U}_{c}}\) . For all \(x\in\mathrm{C}\) such that \(x\neq c\) , we also have \(\overline{\mathrm{U}_{x}}\subset\overline{\mathrm{U}_{c}}\) and \(\overline{\mathrm{U}_{x}}\neq\overline{\mathrm{U}_{c}}\) . By maximality of S, \(\mathrm{C}=\{c\}\) , therefore \(\overline{\mathrm{U}_{c}}=\{c\}\) , which contradicts the hypothesis that \(\mathrm{U}_{c}\) is a nonempty open and closed subset of \(\mathrm{X} - \{c\}\) .

COROLLARY. --- Let X be a compact connected topological space, and let N be the set of points of X whose complement is connected. The set X is the unique compact connected subset of X that contains N.

Let S be a connected and compact subset of X such that \(N \subset S\) . Suppose that \(S \neq X\) and let \(x \in X - S\) . By hypothesis, \(X - \{x\}\) is not connected; there therefore exist open and closed subsets U and V of \(X - \{x\}\) , disjoint and nonempty, such that \(X - \{x\} = U \cup V\) . We may assume that \(S \subset V\) . We have \(\overline{U} = U \cup \{x\}\) and \(\overline{V} = V \cup \{x\}\) , and these spaces are connected (III, p.~278, prop. 17, a)). By assertion c) of this proposition, there exists \(y \in U\) such that \(X - \{y\}\) is connected, which contradicts the inclusions \(N \subset S \subset V\) .

Lemma 7. --- Let T be a locally connected topological space. Let \(\mathcal{U}\) be a directed increasing family of open subsets of T, whose union is T. For \(t \in T\) and \(U \in \mathcal{U}\) , denote by \(U_{t}\) the connected component of t in U if \(t \in U\) and let us set \(U_{t} = \varnothing\) if \(t \notin U\) . For all \(t \in T\) , \(\bigcup_{U \in \mathcal{U}} U_{t}\) is the connected component of t in T.

For \(t \in \mathrm{T}\) , set \(\mathrm{C}_t = \bigcup_{\mathrm{U} \in \mathcal{U}} \mathrm{U}_t\) . We have \(t \in \mathrm{C}_t\) . The sets \(\mathrm{U}_t\) are open and connected, and the same holds for \(\mathrm{C}_t\) because we have \(t \in \mathrm{U}_t\) if \(\mathrm{U}_t \neq \varnothing\) (TG, I, p.~81, prop. 2). Let \(u\) and \(v\) be points of T such that \(\mathrm{C}_u \cap \mathrm{C}_v \neq \varnothing\) . Let \(t\) be a point of \(\mathrm{C}_u \cap \mathrm{C}_v\) ; let \(\mathrm{U} \in \mathcal{U}\) such that \(t \in \mathrm{U}_u\) and let \(\mathrm{V} \in \mathcal{U}\) such that \(v \in \mathrm{V}_v\) . Since \(\mathcal{U}\) is increasingly filtered, there exists \(\mathrm{W} \in \mathcal{U}\) which contains \(\mathrm{U} \cup \mathrm{V}\) . Then, \(\mathrm{W}_u \cap \mathrm{W}_v \neq \varnothing\) , therefore \(\mathrm{W}_u = \mathrm{W}_v\) . More generally, we have \(\mathrm{W}_u' = \mathrm{W}_v'\) for all \(\mathrm{W}' \in \mathcal{U}\) such that \(\mathrm{W} \subset \mathrm{W}'\) , therefore \(\mathrm{C}_u = \mathrm{C}_v\) . This shows that the sets of the form \(\mathrm{C}_t\) form a partition of T into connected open subsets. Consequently, for every \(t \in \mathrm{T}\) , \(\mathrm{C}_t\) is the connected component of \(t\) in T.

THEOREM 2. --- Let X be a compact connected topological space having a countable dense subset. Let a, b be distinct points of X. The following conditions are equivalent:

\begin{enumerate}
\def\labelenumi{(\roman{enumi})}
\item
  There exists a homeomorphism \(f\colon\mathbf{I}\to\mathbf{X}\) such that \(f(0)=a\) and \(f(1)=b\) ;
\item
  Every connected subset of X that contains \(\{a, b\}\) is equal to X;
\item
  The space X is locally connected, and every connected and compact subset of X that contains \(\{a,b\}\) is equal to X;
\item
  For every \(x \in X - \{a, b\}\) , the space \(X - \{x\}\) is not connected.
\end{enumerate}

Assertion (i) implies all the others.

Let \(x\) be a point of \(X - \{a, b\}\) . Since \(X - \{x\}\) contains \(\{a, b\}\) , assertion (ii) implies that this is not a connected subset of \(X\) , so that (ii) implies (iv).

Let us show that (iii) implies (iv). Suppose the hypotheses of (iii) are satisfied. Let \(x \in X - \{a, b\}\) and suppose that \(X - \{x\}\) is connected. Let T be the space \(X - \{x\}\) and let \(\mathcal{U}\) be the set of subsets of T of the form \(X - V\) , where V is a compact neighborhood of \(x\) . By Lemma 7, there exists a compact neighborhood V of \(x\) such that \(a\) and \(b\) belong to the same connected component of \(X - V\) . In particular, they belong to the same connected component of \(X - \mathring{V}\) ; this component is a compact, connected subset of X, distinct from X, which contradicts hypothesis (iii).

It remains to show that assertion (iv) implies condition (i).

Let \(x\) be a point of \(X - \{a, b\}\) and let U, V be open and closed, disjoint and nonempty subsets of \(X - \{x\}\) , such that \(X - \{x\} = U \cup V\) . By III, p.~278, prop. 17, c), there exists a point of U (resp. of V) whose complement in X is connected. Since X admits only two such points, \(a\) and \(b\) , one of them, say \(a\) , is contained in U and the other in V. By loc. cit., a nonempty open and closed subset \(U'\) of U contains a point of \(\{a, b\}\) , hence contains \(a\) . Applying this to \(U - U'\) , it follows that \(U = U'\) . Consequently, U is connected; it is the connected component of \(a\) in \(X - \{x\}\) . Likewise, V is the connected component of \(b\) in \(X - \{x\}\) .

For every \(x \in X - \{a, b\}\) , let \(U_x\) and \(V_x\) denote the connected components of a and b in \(X - \{x\}\) . By the preceding, they are open and closed in \(X - \{x\}\) , disjoint, and their union is equal to \(X - \{x\}\) .

Denote by \(\preccurlyeq\) the relation in X defined as follows: on the one hand, \(a \preccurlyeq x\) and \(x \preccurlyeq b\) for all \(x \in \mathrm{X}\) ; on the other hand, if \(x\) and \(y\) belong to \(\mathrm{X} - \{a, b\}\) , then \(x \preccurlyeq y\) if \(x \in \overline{\mathrm{U}_y}\) . For \(x\) and \(y\) in \(\mathrm{X} - \{a, b\}\) , the subsets \(\mathrm{V}_x\) and \(\mathrm{V}_y\) have the point b in common, and the relation \(x \preccurlyeq y\) is in fact equivalent to the assertion \(\overline{\mathrm{U}_x} \subset \overline{\mathrm{U}_y}\) (III, p.~278, prop. 17, b)). It follows that the relation \(\preccurlyeq\) is an order relation on X.

Let \(x\) and \(y\) be points of X such that one does not have \(x \preccurlyeq y\) . Necessarily, \(x \neq a\) and \(y \neq b\) , and \(x \in V_y\) . If \(x = b\) or \(y = a\) , we have \(y \preccurlyeq x\) . Suppose then that \(x\) and \(y\) are distinct from \(a\) and \(b\) . Since the parts \(U_x\) and \(U_y\) have the point a in common, we have the inclusion \(\overline{V_x} \subset \overline{V_y}\) (loc. cit.), whence, taking the closures of the complements, \(\overline{U_y} \subset \overline{U_x}\) and, a fortiori, \(y \preccurlyeq x\) . The relation \(\preccurlyeq\) in the space X is therefore a total order relation. For every \(x \in X - \{a, b\}\) , we moreover have \(U_x = ] \leftarrow, x[\) and \(V_x = ]x, \rightarrow[\) ; for \(x, y \in X - \{a, b\}\) , we have {]}x, y{[} = \(U_y \cap V_x\) . When x and y range over the points of \(X - \{a, b\}\) , the sets \(U_y \cap V_x\) , the sets \(U_y\) and the sets \(V_x\) form a base for a topology on X. Let \(\widetilde{X}\) denote the corresponding topological space and let i: \(X \to \widetilde{X}\) be the identity map of X. Since, for every \(x \in X\) , the sets \(U_x\) and \(V_x\) are open in X, the map i is continuous.

The space \(\widetilde{\mathbf{X}}\) is separated. Indeed, let \(x\) and \(y\) be distinct points of \(\widetilde{\mathbf{X}}\) such that \(x \preccurlyeq y\) . The sets \(]\leftarrow, y[\) and \(]x, \rightarrow[\) are open neighborhoods of \(x\) and \(y\) ; if they have a common point \(z\) , then \(]\leftarrow, z[\) and \(]z, \rightarrow[\) are disjoint open neighborhoods of \(x\) and \(y\) . Since X is compact, the map \(i\) is therefore a homeomorphism (TG, I, p.~63, cor. 2).

Therefore, the image under \(i\) of a countable dense subset of X intersects every nonempty open interval of the ordered set \((\mathrm{X}, \preccurlyeq)\) . It then follows from prop. 16 of III, p.~276 that there exists an isomorphism \(c\) of the ordered set \(\mathbf{I}\) onto \((\mathrm{X}, \preccurlyeq)\) . According to the remark following this proposition, this isomorphism is a homeomorphism of \(\mathbf{I}\) onto the topological space \(\widetilde{\mathrm{X}}\) . The map \(f = i^{-1} \circ c\) is then a homeomorphism of \(\mathbf{I}\) onto X which maps 0 to \(a\) and 1 to \(b\) .

\subsection{Injective paths}

PROPOSITION 18. --- Let X be a Hausdorff topological space. Let a and b be distinct points of X that belong to the same path-connected component of X. There exists an injective path connecting a to b in X.

Let \(f \colon \mathbf{I} \to \mathbf{X}\) be a path connecting \(a\) to \(b\) in \(\mathbf{X}\) . Let \(\mathcal{U}\) be the set of open subsets \(\mathrm{U}\) of {]}0,1{[} such that \(f(x) = f(y)\) for every connected component {]}x, y{[} of U.

Lemma 8. --- The set U, ordered by inclusion, is inductive.

Let V be a totally ordered subset of U. Let us prove that the union V of the sets belonging to V is an element of U, that is, that, for every connected component {]}u, v{[} of V, we have \(f(u) = f(v)\) .

Let \(x\) be a point of \([u, v[\) . Let \(\mathcal{V}_x\) be the set of \(U \in \mathcal{V}\) such that \(x \in U\) ; for such a \(U\) , denote by \([u_U, v_U[\) the connected component of \(x\) in \(U\) . We have \(u_{U'} \leqslant u_U < v_U \leqslant v_{U'}\) if \(U\) and \(U'\) are elements of \(\mathcal{V}_x\) such that \(U \subset U'\) . Since the union over \(U \in \mathcal{V}_x\) of the \([u_U, v_U[\) is equal to {]}u, v{[} by Lemma 7 of III, p.~280, we have u = lim \(u_{U}\) and v = lim \(v_{U}\) , where the limits are taken over the directed set \(V_{x}\) . Since the map f is continuous and the topological space X is Hausdorff, the equalities \(f(u_{\mathrm{U}}) = f(v_{\mathrm{U}})\) for all \(U \in V_{x}\) imply that \(f(u) = f(v)\) , which is what had to be proved.

By virtue of E, III, p.~20, th. 2, there exists an open subset U belonging to \(\mathcal{U}\) which is maximal for the inclusion relation.

Let \(g \colon \mathbf{I} \to \mathrm{X}\) be the map defined in the following way. If \(t \notin \mathrm{U}\) we set \(g(t) = f(t)\) ; otherwise, denoting by \(]u, v[\) the connected component of \(t\) in \(\mathrm{U}\) we set \(g(t) = f(u) = f(v)\) , so that the map \(g|[u, v]\) is constant with image \(f(u)\) .

Prop. 18 follows from the following lemma.

Lemma 9. --- There exists a continuous, increasing, and surjective map \(u\colon\mathbf{I}\to\mathbf{I}\) and an injective path \(c\colon\mathbf{I}\to\mathbf{X}\) connecting \(a\) to \(b\) such that \(g=c\circ u\) .

Let us first prove that the map \(g\) is continuous. Let \(t\) be a point of \(\mathbf{I}\) . If \(t \in \mathrm{U}\) , \(g\) is constant in a neighborhood of \(t\) , hence continuous at \(t\) . Suppose \(t \notin \mathrm{U}\) and let \(\mathrm{W}\) be an open neighborhood of \(g(t)\) in \(\mathrm{X}\) . Let \(\mathrm{V}\) be an open interval in \(\mathbf{I}\) containing \(t\) such that \(f(\mathrm{V}) \subset \mathrm{W}\) ; such an interval exists since \(f\) is continuous at \(t\) . Let us prove that \(g(\mathrm{V}) \subset \mathrm{W}\) . Let \(x \in \mathrm{V}\) ; if \(x \notin \mathrm{U}\) , we have \(g(x) = f(x)\) , therefore \(g(x) \in \mathrm{W}\) . Suppose then \(x \in \mathrm{U}\) and let \(]u, v[\) be the connected component of \(x\) in \(\mathrm{U}\) . Observe that \(]u, v[\) does not contain \(t\) . Consequently, if \(x < t\) , we have \(x < v \leqslant t\) , whence \(v \in \mathrm{V}\) and \(g(x) = g(v) = f(v) \in f(\mathrm{V}) \subset \mathrm{W}\) . One proves similarly that \(g(x) \in \mathrm{W}\) if \(x > t\) . Thus, \(g\) is continuous at \(t\) .

Let \(u\) and \(v\) be points of \(\mathbf{I}\) such that \(g(u) = g(v)\) and \(u < v\) . Let us prove that \(]u, v[ \subset U\) . Set \(U' = U \cup ]u, v[\) ; it is an open subset of \(]0, 1[\) .

If \(u\) belongs to U, denote by \(u'\) the greatest lower bound in I of the connected component of \(u\) in U; set \(u' = u\) if \(u\) does not belong to U. Likewise, if \(v\) belongs to U, denote by \(v'\) the least upper bound in I of the connected component of \(v\) in U; set \(v' = v\) if \(v\) does not belong to U. Then, \(]u', v'[\) is a connected component of U', and we have \(f(u') = g(u) = g(v) = f(v')\) . Since the connected components of U' distinct from \(]u', v'[\) are connected components of U, the open set U' is an element of \(\mathcal{U}\) . Since U is a maximal element of \(\mathcal{U}\) for the inclusion relation, we have U' = U and \(]u, v[\) is contained in U. This proves that \(g\) is constant on the interval \([u, v]\) . The fibers of \(g\) are therefore intervals of I, and these intervals are closed because g is continuous.

Let R denote the equivalence relation associated with g, and let p be the canonical surjection from I onto I/R. There exists a unique continuous map \(g'\) from I/R to X such that \(g = g' \circ p\) ; this map is injective. By the corollary, III, p.~276, of Prop. 16, there exists a map u, increasing, continuous, and surjective, such that R is the equivalence relation associated with u. Since the space I is compact and the space X is Hausdorff, the map u is closed, hence strict (I, p.~18, Example 2). It defines, by passing to the quotient, a homeomorphism \(u'\) from I/R onto I. Then, the map \(g' \circ (u')^{-1}\) from I into X is an injective path with initial point a and terminal point b.

\subsection{Lifting of paths}

THEOREM 3. --- Let I be an interval of R and let X be a nonempty separated topological space. Let p: X → I be a continuous, open, and proper map whose fibers are totally disconnected (TG, I, p.~83). The map p is surjective. For every point x of X, there exists a continuous section s of p such that s(p(x)) = x.

The set \(p(\mathrm{X})\) is an open, closed (TG, I, p.~72, prop. 1), nonempty subset of I. Since I is connected, we have \(p(\mathrm{X}) = \mathrm{I}\) ; the map \(p\) is therefore surjective.

For every pair \((a,b)\in\mathrm{I}\times\mathrm{I}\) such that \(a\leqslant b\) , denote by \(F_{a,b}\) the set of pairs \((y,z)\in\overline{p}^{-1}(a)\times\overline{p}^{-1}(b)\) such that y and z belong to the same connected component of \(\overline{p}^{-1}([a,b])\) .

Lemma 10. --- Let \(a\) and \(b\) be points of I such that \(a \leqslant b\) .

\begin{enumerate}
\def\labelenumi{\alph{enumi})}
\item
  \(F_{a,b}\) is closed in \(\overline{p}^{-1}(a)\times\overline{p}^{-1}(b)\).
\item
  We have \(\operatorname{pr}_1(\mathrm{F}_{a,b}) = \overline{p}^{-1}(a)\) and \(\operatorname{pr}_2(\mathrm{F}_{a,b}) = \overline{p}^{-1}(b)\) .
\item
  Let \(c \in I\) such that \(b \leqslant c\) . If \((y,z)\) belongs to \(F_{a,b}\) and \((z,t)\) belongs to \(F_{b,c}\) then \((y,t)\) belongs to \(F_{a,c}\) .
\end{enumerate}

The set \(\overline{p}^{-1}([a,b])\) is compact (TG, I, p.~77, prop. 7). Consequently, for a pair \((y,z)\in\overline{p}^{-1}(a)\times\overline{p}^{-1}(b)\) to belong to \(F_{a,b}\) , it is necessary and sufficient that every open and closed subset of \(\overline{p}^{-1}([a,b])\) which contains \(y\) contains \(z\) (TG, II, p.~32, prop. 6).

Let us set \(\mathrm{Y} = \overline{p}^{-1}([a,b])\) . For every subset U of Y that is both open and closed, the set \(((\mathrm{U} \times \mathrm{U}) \cup ((\mathrm{Y} - \mathrm{U}) \times (\mathrm{Y} - \mathrm{U}))) \cap (\overline{p}^{-1}(a) \times \overline{p}^{-1}(b))\) is closed in \(\overline{p}^{-1}(a) \times \overline{p}^{-1}(b)\) . The intersection of these sets is equal to \(\mathrm{F}_{a,b}\) , whence \(a\) ).

Let \(y \in \overline{p}^{-1}(a)\) . Let \(\mathcal{U}\) denote the set of open and closed neighborhoods of y in Y. The map from Y to {[}a, b{]} deduced from p by passing to subspaces is open and proper (I, p.~17, prop. 8), hence also closed. It follows that, for every \(U\) belonging to \(\mathcal{U}\) , \(p(U)\) is a nonempty open and closed subset of {[}a, b{]}; since the interval {[}a, b{]} is connected, one has \(p(U) = [a, b]\) and in particular \(U \cap \overline{p}^{-1}(b) \neq \varnothing\) . Consequently, \((U \cap \overline{p}^{-1}(b))_{U \in \mathcal{U}}\) is a decreasing filtered family of nonempty closed subsets of the compact space \(\overline{p}^{-1}(b)\) . The intersection of this family is nonempty (TG, I, p.~59); let z be one of its elements. One has \((y, z) \in \mathrm{F}_{a,b}\) . We have proved the relation \(\mathrm{pr}_{1}(\mathrm{F}_{a,b}) = \overline{p}^{-1}(a)\) . The relation \(\mathrm{pr}_{2}(\mathrm{F}_{a,b}) = \overline{p}^{-1}(b)\) is deduced from it by replacing p, I, a, b by -p, -I, -b, -a.

Under the hypotheses of \(c\) ), the pair \((y,t)\) belongs to \(\overline{p}^{-1}(a)\times\overline{p}^{-1}(c)\) . The set \(\{y,z\}\) is contained in a connected subset C of \(\overline{p}^{-1}([a,b])\) and the set \(\{z,t\}\) is contained in a connected subset C' of \(\overline{p}^{-1}([b,c])\) . Then, C \(\cup\) C' is a connected subset of \(\overline{p}^{-1}([a,c])\) (TG, I, p.~81, prop. 2) and contains \(\{y,t\}\) , whence the relation \((y,t)\in\mathrm{F}_{a,c}\) .

Returning to the proof of Th. 3. Each fiber of the map p is compact (TG, I, p.~77, prop. 7). By Tychonoff's theorem (TG, I, p.~63, th. 3), the product space \(K = \prod_{a \in I} \overline{p}^{-1}(a)\) is compact. Let \(K'\) be the set of elements \((y_{a})_{a \in I}\) of K such that \(y_{p(x)}\) equals x and that one has \((y_{a}, y_{b}) \in F_{a,b}\) for every pair \((a, b)\) of elements of I such that a \textless{} b. Theorem 3 follows from the following lemma.

Lemma 11. --- a) The set K' is not empty.

\begin{enumerate}
\def\labelenumi{\alph{enumi})}
\setcounter{enumi}{1}
\tightlist
\item
  Let \((s_{a})_{a\in\mathrm{I}}\) be an element of \(K'\) . The map \(s\colon a\mapsto s_{a}\) from I to X is a continuous section of p such that \(s(p(x))=x\) .
\end{enumerate}

For every finite subset S of I containing the point \(p(x)\) , denote by \(K_{S}\) the set of elements \((y_{a})_{a\in\mathrm{I}}\) of K satisfying the relation \(y_{p(x)} = x\) and the relations \((y_{a}, y_{b}) \in \mathrm{F}_{a,b}\) for every pair \((a, b)\) of elements of S such that a \textless{} b. The sets \(K_{S}\) are closed in K (Lemma 10, a)) and form a decreasing filtered family of subsets of K, whose intersection \(K'\) .

To prove that \(K'\) is nonempty, it suffices to show that, for every finite subset S of I containing the point \(p(x)\) , the set \(K_{S}\) is not empty (TG, I, p.~59).

Let S be such a subset; arrange its elements in a strictly increasing sequence \((a_{1},\ldots,a_{n})\), and let i denote the integer such that \(p(x)=a_{i}\) . Set \(y_{a_{i}}=x\) . Lemma 10, b), allows one to construct by induction elements \(y_{a_{j}}\in\overline{p}^{-1}(a_{j})\) , for \(i<j\leqslant n\) , and, by descending induction, elements \(y_{a_{j}}\in\overline{p}^{-1}(a_{j})\) for \(1\leqslant j<i\) , so that one has \((y_{a_{j}},y_{a_{j+1}})\in\mathrm{F}_{a_{j},a_{j+1}}\) for every integer j such that \(1\leqslant j<n\) . By Lemma 10, c), we have \((y_{a},y_{b})\in\mathrm{F}_{a,b}\) for every pair \((a,b)\) of elements of S such that a\textless b. Since the map p is surjective, we can choose for every \(a\in I-S\) an element \(y_{a}\in\overline{p}^{-1}(a)\) . The family \((y_{a})_{a\in I}\) thus constructed belongs to \(K_{S}\) , therefore \(K_{S}\) is not empty.

Let us prove \(b\) ). By definition of \(K'\) , \(s\) is a section of \(p\) such that \(s(p(x)) = x\) . Let \(a \in I\) . Let us prove the continuity of \(s\) at the point \(a\) . Let U be an open neighborhood of \(s_a\) in X. Since \(\overline{p}^{-1}(a)\) is a compact space (TG, I, p.~77, prop. 7) and totally disconnected, \(s_a\) has, in \(\overline{p}^{-1}(a)\) , an open and closed neighborhood C contained in U (TG, II, p.~32, corollary). The sets C and \(\overline{p}^{-1}(a) - C\) are closed in \(\overline{p}^{-1}(a)\) ; therefore compact, and they are disjoint. Since X is separated, they have disjoint open neighborhoods V and \(V'\) . The set \((V \cap U) \cup V'\) is an open neighborhood of the fiber \(\overline{p}^{-1}(a)\) in X; since the map \(p\) is closed (TG, I, p.~72, prop. 1), \((V \cap U) \cup V'\) contains a set of the form \(\overline{p}^{-1}(J)\) , where \(J \subset I\) is an open interval containing \(a\) (I, p.~75, lemma). Put \(W = V \cap U \cap \overline{p}^{-1}(J)\) . The set W is open in \(\overline{p}^{-1}(J)\) ; it is also closed in \(\overline{p}^{-1}(J)\) since we have \(\overline{p}^{-1}(J) - W = V' \cap \overline{p}^{-1}(J)\) . Let \(b \in J\) . The closed interval of I with endpoints \(a\) and \(b\) is contained in J. By hypothesis, \((s_a, s_b)\) belongs to \(F_{a,b}\) if \(a \leqslant b\) and \((s_b, s_a)\) belongs to \(F_{b,a}\) if \(b \leqslant a\) . There therefore exists a connected subset of \(\overline{p}^{-1}(J)\) containing \(\{s_a, s_b\}\) . Consequently, the point \(s_b\) belongs to every open and closed subset of \(\overline{p}^{-1}(J)\) which contains \(s_a\) , hence in particular to W; a fortiori, \(s_b\) belongs to U. We therefore have \(s(J) \subset U\) , which proves the continuity of \(s\) at the point \(a\) .

COROLLARY. --- Let X and B be topological spaces and \(p: X \to B\) a continuous, open, proper and separated map whose fibers are totally disconnected. Let I be an interval of R, let \(f: I \to B\) be a continuous map, let a be a point of I and let x be a point of X such that \(f(a) = p(x)\) . There exists a continuous map \(g: I \to X\) such that \(p \circ g = f\) and \(g(a) = x\) .

Let us set \(X' = I \times_{B} X\) and denote by \(p'\colon X'\to I\) and \(f'\colon X'\to X\) the canonical projections. The map \(p'\) is continuous, open, proper and separated (I, p.~17, prop. 8 and p.~27, prop. 4). Since the space I is separated, the space \(X'\) is separated (I, p.~26, remark 3). The fibers of \(p'\) are totally disconnected (I, p.~10, corollary, a)). By th. 3, there exists a continuous section \(s'\) of \(p'\) which takes the value \((a,x)\) at \(a\) . The map \(g = f'\circ s'\) from I to X is continuous; we have \(p\circ g = p\circ f'\circ s' = f\circ p'\circ s' = f\) and \(g(a) = x\) , whence the corollary.

THEOREM 4. --- Let X be a topological space, let G be a discrete group acting properly on X and let p: X → X/G be the canonical map. Let I be an interval of R, let f: I → X/G be a continuous map, let a be a point of I and let x be a point of X such that \(f(a) = p(x)\) . There exists a continuous map \(\varphi\) : I → X such that \(p \circ \varphi = f\) and \(\varphi(a) = x\) .

Let us first treat the case where I is a closed bounded interval of R.

By TG, III, p.~29, prop. 3, the space X is separated.

Let \(y\) be a point of X/G and let \(x \in X\) such that \(y = p(x)\) . The stabilizer \(K_x\) of \(x\) is a finite subgroup of G; moreover, there exist neighborhoods \(U_x\) of \(x\) in X and \(V_y\) of \(y\) in X/G such that \(U_x\) is stable under \(K_x\) , \(gU_x \cap U_x = \varnothing\) if \(g \notin K_x\) , and \(p\) induces a homeomorphism from \(U_x/K_x\) onto \(V_y\) . In addition, for every \(g \in G\) , \(p\) induces a homeomorphism from \(gU_x\) onto \(V_y\) . Since I is compact, there exist integers \(m\) and \(n\) such that \(m \leqslant 0 \leqslant n\) and a finite sequence \((a_i)_{m \leqslant i \leqslant n}\) of elements of I such that \(a_0 = a\) , \(I = {[}a_m,a_n{]}\) , and such that, for every \(i \in \{m, \ldots, n-1\}\) , \(f({[}a_i,a_{i+1}{]})\) is contained in an open set \(V_{y_i}\) of X/G constructed as above.

Let \(x_0\) be the unique element of \(\overline{p}^{-1}(y_0)\) such that \(x \in \mathrm{U}_{x_0}\) . Let \(q_0\) be the canonical map from \(\mathrm{U}_{x_0}\) onto \(\mathrm{U}_{x_0}/\mathrm{K}_{x_0}\) by passing to the quotient. The map \(p\) induces a homeomorphism \(i_0\) of \(\mathrm{U}_{x_0}/\mathrm{K}_{x_0}\) onto \(\mathrm{V}_{y_0}\) such that \(i_0 \circ q_0 = p \mid \mathrm{U}_{x_0}\) . The map \(q_0\) is proper (TG, III, p.~29, prop. 3), open (TG, I, p.~31, example 1), and separated, since X is separated. Its fibers are totally disconnected, since they are finite. By the corollary, III, p.~286, there exists a continuous map \(\varphi_0\colon [a_0,a_1]\to \mathrm{U}_{x_0}\) such that \(p\circ \varphi_0 = f\mid [a_0,a_1]\) .

We construct similarly, by induction on the integer \(i \in \{0, \ldots, n-1\}\) , a point \(x_i \in \overline{p}^{-1}(y_i)\) and a continuous map \(\varphi_i: [a_i, a_{i+1}] \to X\) whose image is contained in \(\mathrm{U}_{x_i}\) such that \(p \circ \varphi_i = f \mid [a_i, a_{i+1}]\) and such that \(\varphi_i(a_{i+1}) = \varphi_{i+1}(a_{i+1})\) if \(0 \leqslant i < n-1\) .

Analogously, one constructs by descending induction on the integer \(i \in \{m, \ldots, -1\}\) , a point \(x_i \in \overline{p}^{-1}(y_i)\) and a continuous map \(\varphi_i: [a_i, a_{i+1}] \to X\) whose image is contained in \(U_{x_i}\) such that \(p \circ \varphi_i = f \mid [a_i, a_{i+1}]\) and such that \(\varphi_i(a_{i+1}) = \varphi_{i+1}(a_{i+1})\) if \(m \leqslant i < 0\) .

There exists a unique map \(\varphi\colon\mathrm{I}\to\mathrm{X}\) which coincides with \(\varphi_{i}\) on \([a_{i},a_{i+1}]\) for \(m\leqslant i<n\) . It is continuous (TG, I, p.~19, prop. 4). It is a continuous lifting of \(f\) to X such that \(\varphi(a)=x\) .

This proves the theorem when I is compact. In the general case, there exist sequences \((a_{n})_{n\in\mathbf{N}}\) and \((b_{n})_{n\in\mathbf{N}}\) such that \((a_{n})\) is stationary with limit \(\inf(I)\), \((b_{n})\) is stationary with limit \(\sup(I)\), \(a=a_{0}=b_{0}\) , and such that \((a_{n})\) (resp. \((b_{n})\) ) is constant if I has a least (resp. a greatest) element. By the preceding, for every integer \(n\in N\) there exists a continuous lifting \(\varphi_{n}\) of \(f\mid[a_{n},a_{n+1}]\) to X and a continuous lifting \(\varphi_{n}^{\prime}\) of \(f\mid[b_{n+1},b_n]\) to X such that \(\varphi_{0}(a_{0})=\varphi_{0}^{\prime}(b_{0})=x\) and \(\varphi_{n+1}(a_{n+1})=\varphi_{n}(a_{n+1})\), \(\varphi_{n+1}^{\prime}(b_{n+1})=\varphi_{n}^{\prime}(b_{n+1})\) . There exists a unique map \(\varphi\colon I\to X\) which coincides with \(\varphi_{n}\) on \([a_{n},a_{n+1}]\) and with \(\varphi_{n}^{\prime}\) on \([b_{n+1},b_{n}]\) , for every \(n\in N\) . It is continuous (loc. cit.) and it is a continuous lifting of f to X such that \(\varphi(a)=x\) . The theorem is thus proved.

Examples. --- 1) Let X be a separated topological space and let G be a finite group, endowed with the discrete topology, acting continuously on X. The action is then proper (TG, III, p.~28, prop. 2). The assertion of Theorem 4 then follows directly from the corollary of Theorem 3.

\begin{enumerate}
\def\labelenumi{\arabic{enumi})}
\setcounter{enumi}{1}
\tightlist
\item
  Let n be an integer \(\geqslant 0\) . Denote by \(P_{n}\) the set of monic polynomials \(P \in C[X]\) of degree n, endowed with the topology for which the map \((c_{0}, \ldots, c_{n-1}) \mapsto X^{n} + c_{n-1}X^{n-1} + \cdots + c_{0}\) is a homeomorphism of \(C^{n}\) onto \(P_{n}\) . The map p from \(C^{n}\) to \(P_{n}\) defined by \(p(z_{1}, \ldots, z_{n}) = (X - z_{1}) \ldots (X - z_{n})\) is continuous. The symmetric group \(S_{n}\) acts on \(C^{n}\) by permutation of the factors and p defines, by passage to the quotient, a homeomorphism from \(C^{n}/S_{n}\) onto \(P_{n}\) (TG, VIII, p.~22, prop. 1, I, p.~23, cor. 1 and TG, VIII, p.~20). We therefore deduce the following statement:
\end{enumerate}

Let I be an interval of R, let \((c_{0},\ldots,c_{n-1})\) be a sequence of continuous maps from I to C, let a be a point of I and let \((z_{1},\ldots,z_{n})\) be a sequence of complex numbers such that \((\mathrm{X}-z_{1})\ldots(\mathrm{X}-z_{n})=\mathrm{X}^{n}+c_{n-1}(a)\mathrm{X}^{n-1}+\cdots+c_{0}(a)\) . There exists a sequence \((\lambda_{1},\ldots,\lambda_{n})\) of continuous maps from I to C such that \(\lambda_{i}(a)=z_{i}\) for \(1\leqslant i\leqslant n\) and \((\mathrm{X}-\lambda_{1}(t))\ldots(\mathrm{X}-\lambda_{n}(t))=\mathrm{X}^{n}+c_{n-1}(t)\mathrm{X}^{n-1}+\cdots+c_{0}(t)\) for all \(t\in I\) .

\section{POINCARÉ GROUPOID}

\subsection{Poincaré groupoid}

DEFINITION 1. --- Let X be a topological space, let \(c_{0}\) and \(c_{1}\) be paths in X and let \(\sigma\colon\mathbf{I}\times\mathbf{I}\to\mathbf{X}\) be a homotopy connecting \(c_{0}\) to \(c_{1}\) . One says that \(\sigma\) is a strict homotopy if the maps \(s\mapsto\sigma(0,s)\) and \(s\mapsto\sigma(1,s)\) are constant.

We say that paths \(c_{0}\) and \(c_{1}\) in X are strictly homotopic if there exists a strict homotopy connecting \(c_{0}\) to \(c_{1}\) .

Two strictly homotopic paths have the same origin and the same endpoint.

Example. --- Let c be a path in X and let \(\varphi\colon I\to I\) be a continuous map such that \(\varphi(0)=0\) and \(\varphi(1)=1\) . The paths c and \(c\circ\varphi\) are strictly homotopic. Indeed, the map \(\sigma\colon I\times I\to X\) defined by \(\sigma(t,s)=c((1-s)t+s\varphi(t))\) is a strict homotopy connecting c to \(c\circ\varphi\) .

Let X be a topological space. Recall (cf.~III, p.~257) that \(\Lambda(X)\) denotes the topological space \(\mathcal{C}_c(\mathbf{I};\mathrm{X})\) of paths in X and that for \(x\) , \(y \in X\) , \(\Lambda_{x,y}(X)\) is the subspace of \(\Lambda(X)\) formed by the paths with origin \(x\) and endpoint \(y\) . The family of sets \(\Lambda_{x,y}(X)\) , for \(x\) , \(y \in X\) , is a partition of the space of paths of X. By the canonical bijection (III, p.~257, remark 2) of \(\mathcal{C}(\mathbf{I} \times \mathbf{I};\mathrm{X})\) onto \(\mathcal{C}(\mathbf{I};\Lambda(\mathrm{X}))\) , strict homotopies correspond to paths \(c: \mathbf{I} \to \Lambda(X)\) whose image is contained in a subspace of the form \(\Lambda_{x,y}(\mathrm{X})\) . The relation “the paths \(c_0\) and \(c_1\) are strictly homotopic” is therefore an equivalence relation in \(\Lambda(\mathrm{X})\) (III, p.~259, prop. 5) and the equivalence classes for this relation are the arcwise connected components of the subspaces \(\Lambda_{x,y}(\mathrm{X})\) of the space of paths of X. We denote by \(\varpi_{x,y}(\mathrm{X})\) the set \(\pi_0(\Lambda_{x,y}(\mathrm{X}))\) and call a class of paths joining \(x\) to \(y\) any element of \(\varpi_{x,y}(\mathrm{X})\) .

DEFINITION 2. --- We call the loop space of X, and denote by \(\Omega(X)\) , the subspace of \(\Lambda(X)\) consisting of the loops (III, p.~256, def. 1) in X.

We denote \(\Omega_x(\mathrm{X})\) the set \(\Lambda_{x,x}(\mathrm{X})\) . The elements of \(\Omega_x(\mathrm{X})\) are called the loops in X at \(x\) and the elements of \(\varpi_{x,x}(\mathrm{X})\) are called loop classes in X at \(x\) . The map \(e_x: \mathbf{I} \to \mathbf{X}\) is constant with image \(x\) and is a loop, called the constant loop at \(x\) ; its strict homotopy class is denoted by \(\varepsilon_x\) . The map \(x \mapsto e_x\) from X to \(\Lambda(\mathrm{X})\) is continuous (III, p.~257, prop. 1).

Let X be a topological space and let x, y, z be points of X. By passing to arcwise connected components, we deduce from the continuous map \(c \mapsto \overline{c}\) from \(\Lambda_{x,y}(X)\) to \(\Lambda_{y,x}(X)\) (III, p.~258, corollary) a map from \(\varpi_{x,y}(X)\) to \(\varpi_{y,x}(X)\) , denoted \(\gamma \mapsto \overline{\gamma}\) . If \(\gamma \in \varpi_{x,y}(X)\) , \(\overline{\gamma}\) is called the inverse of the path class \(\gamma\) .

Similarly, if we identify the sets \(\pi_0(\Lambda_{x,y}(\mathrm{X})) \times \pi_0(\Lambda_{y,z}(\mathrm{X}))\) and \(\pi_0(\Lambda_{x,y}(\mathrm{X}) \times \Lambda_{y,z}(\mathrm{X}))\) (III, p.~260, prop. 6), we deduce from the continuous map \((c,d) \mapsto c * d\) from \(\Lambda_{x,y}(\mathrm{X}) \times \Lambda_{y,z}(\mathrm{X})\) to \(\Lambda_{x,z}(\mathrm{X})\) (III, p.~258, corollary), by passing to arcwise connected components, a map \(C_{x,y,z} \colon \varpi_{x,y}(\mathrm{X}) \times \varpi_{y,z}(\mathrm{X}) \to \varpi_{x,z}(\mathrm{X})\) . For \(\gamma \in \varpi_{x,y}(\mathrm{X})\) and \(\delta \in \varpi_{y,z}(\mathrm{X})\) , write \(\gamma \delta\) for the strict homotopy class \(C_{x,y,z}(\gamma,\delta)\) . It is called the composite of juxtaposable path classes \(\gamma\) and \(\delta\) .

We have \(\overline{\overline{\gamma}}=\gamma\) and \(\overline{\gamma\delta}=\overline{\delta}\overline{\gamma}\) .

PROPOSITION 1. --- Let X be a topological space, let x, y, z, u be points of X, and let \(\gamma_{1} \in \varpi_{x,y}(X)\) , \(\gamma_{2} \in \varpi_{y,z}(X)\) , \(\gamma_{3} \in \varpi_{z,u}(X)\) be path classes. We have

\begin{enumerate}
\def\labelenumi{(\arabic{enumi})}
\tightlist
\item
\end{enumerate}

\[
\varepsilon_ {x} \gamma_ {1} = \gamma_ {1} \varepsilon_ {y} = \gamma_ {1},\tag{2}
\]

\[
\gamma_ {1} \overline {{\gamma}} _ {1} = \varepsilon_ {x}, \qquad \overline {{\gamma}} _ {1} \gamma_ {1} = \varepsilon_ {y},\tag{3}
\]

\[
(\gamma_ {1} \gamma_ {2}) \gamma_ {3} = \gamma_ {1} (\gamma_ {2} \gamma_ {3}).
\]

Let \(c_{1} \in \Lambda_{x,y}(\mathrm{X})\) , \(c_{2} \in \Lambda_{y,z}(\mathrm{X})\) , \(c_{3} \in \Lambda_{z,u}(\mathrm{X})\) be representatives of \(\gamma_{1}\) , \(\gamma_{2}\) and \(\gamma_{3}\) respectively. Let \(\varphi: \mathbf{I} \to \mathbf{I}\) be the function defined by

\[
\varphi (t) = \left\{ \begin{array}{l l} t / 2 & \text { for } 0 \leqslant t \leqslant 1 / 2, \\ t - 1 / 4 & \text { for } 1 / 2 \leqslant t \leqslant 3 / 4, \\ 2 t - 1 & \text { for } 3 / 4 \leqslant t \leqslant 1. \end{array} \right.\tag{4}
\]

The function \(\varphi\) is affine on each of the three intervals \([0,1/2]\) , \([1/2,3/4]\) and \([3/4,1]\) ; it is therefore continuous. We have \(\varphi(0) = 0\) and \(\varphi(1) = 1\) . It follows from formula (1) of III, p.~256 defining the juxtaposition of paths that

\[
c _ {1} * (c _ {2} * c _ {3}) = ((c _ {1} * c _ {2}) * c _ {3}) \circ \varphi .
\]

According to the example in III, p.~289, the paths \(c_{1} * (c_{2} * c_{3})\) and \((c_{1} * c_{2}) * c_{3}\) are strictly homotopic, whence equality (3).

Similarly, the function \(\psi\colon\mathbf{I}\to\mathbf{I}\) defined by

\[
\psi (t) = \left\{ \begin{array}{l l} 2 t & \text { for } 0 \leqslant t \leqslant 1 / 2, \\ 1 & \text { for } 1 / 2 \leqslant t \leqslant 1 \end{array} \right.\tag{5}
\]

is continuous and satisfies \(\psi(0)=0,\psi(1)=1\) . The equality \(\gamma_{1}\varepsilon_{y}=\gamma_{1}\) then follows from the fact that

\[
c _ {1} * e _ {y} = c _ {1} \circ \psi
\]

and the equality \(\varepsilon_{x}\gamma_{1}=\gamma_{1}\) is proved similarly, whence (1).

The map \(\sigma\colon\mathbf{I}\times\mathbf{I}\to\mathrm{X}\) defined by

\[
\sigma (t, s) = \left\{ \begin{array}{l l} c _ {1} (2 t s) & \text { for } 0 \leqslant t \leqslant 1 / 2, \\ c _ {1} (2 (1 - t) s) & \text { for } 1 / 2 \leqslant t \leqslant 1 \end{array} \right.\tag{6}
\]

is continuous; it is a strict homotopy connecting the path \(e_{x}\) to the path \(c_{1}*\overline{c_{1}}\) , whence the first equality in (2). The second follows from the first and from the fact that, for every path c, one has \(c=\overline{\overline{c}}\) .

Remark 1. --- Let X be a topological space. Let n be an integer \(\geqslant 1\) and let \((c_{1},\ldots,c_{n})\) be a sequence of paths in X such that \(c_{i}\) and \(c_{i+1}\) are juxtaposable for \(1 \leqslant i \leqslant n-1\) (such a sequence is called a sequence of juxtaposable paths). Denote by c the path

\[
c _ {1} * \left(c _ {2} * \left(\dots * \left(c _ {n - 1} * c _ {n}\right) \ldots\right)\right)
\]

and \(c'\) the path defined by \(c'(t) = c_i(nt - i + 1)\) for \(1 \leqslant i \leqslant n\) and \(t \in [\frac{i-1}{n}, \frac{i}{n}]\) . The paths \(c\) and \(c'\) have the same image and are strictly homotopic: one is the composite of the other with a homeomorphism of I fixing 0 and 1 (cf. III, p. 289, example). We will sometimes denote \(c_{1} * c_{2} * \cdots * c_{n}\) the path \(c'\) .

There exists a unique directed graph \(\varpi(\mathrm{X})\) whose set of vertices is X and whose set of arrows connecting a point \(x\) to a point \(y\) is \(\varpi_{x,y}(\mathrm{X})\) , and in which the maps \(C_{x,y,z}\) define a composition law. By Proposition 1, \(\varpi(\mathrm{X})\) is a groupoid (II, p.~162, def. 4). For every \(x \in \mathrm{X}\) , the identity element at the vertex \(x\) of this groupoid is the class of the constant loop with image \(x\) . The inverse of an arrow \(\gamma\) is the arrow \(\overline{\gamma}\) , which we shall also denote by \(\gamma^{-1}\) . In particular, the composition law \(C_{x,x,x}\) equips, for each \(x \in \mathrm{X}\) , the set \(\varpi_{x,x}(\mathrm{X})\) with a group structure; we denote this group \(\pi_1(\mathrm{X}, x)\) .

DEFINITION 3. --- Let X be a topological space. The groupoid \(\varpi(X)\) is called the Poincaré groupoid, or fundamental groupoid, of the space X. Let x be a point of X; the group \(\pi_1(X,x)\) of loop classes at x is called the Poincaré group, or fundamental group, of the space X at the point x.

Let \(\mathcal{U}\) be a set of subsets of X whose interiors cover X. The path classes in X whose image is contained in one of the subsets belonging to \(\mathcal{U}\) generate the groupoid \(\varpi(\mathrm{X})\) (lemma 4 of III, p.~272).

The orbits of the groupoid \(\varpi(X)\) (II, p.~162) coincide with the arcwise connected components of the space X. In particular (loc. cit.), one has:

PROPOSITION 2. --- Let X be a topological space.

\begin{enumerate}
\def\labelenumi{\alph{enumi})}
\item
  If the space X is path-connected, the groups \(\pi_1(\mathrm{X},x)\) and \(\pi_1(\mathrm{X},y)\) are isomorphic for all points \(x\) and \(y\) of \(\mathrm{X}\) .
\item
  Let x be a point of X; the following conditions are equivalent:
\end{enumerate}

\begin{enumerate}
\def\labelenumi{(\roman{enumi})}
\item
  The group \(\pi_1(X,x)\) is trivial;
\item
  Two paths with origin x in X that have the same endpoint are strictly homotopic;
\item
  Every loop with origin x in X is strictly homotopic to the constant loop with image x.
\end{enumerate}

More precisely, let X be a path-connected topological space and let x and y be points of X. For every element \(\delta\) of \(\varpi_{x,y}(X)\) , the map \(u_{\delta} \colon \gamma \mapsto \delta \gamma \delta^{-1}\) from \(\pi_1(\mathrm{X},y)\) to \(\pi_1(\mathrm{X},x)\) is a group isomorphism whose inverse isomorphism is \(u_{\delta^{-1}}\) . For \(\delta \in \varpi_{x,y}(\mathrm{X})\) and \(\gamma \in \pi_1(\mathrm{X},y)\) we set \(^{\delta}\gamma = u_{\delta}(\gamma)\) . When \(x = y\) , we have \(^{\delta}\gamma = \operatorname{Int}(\delta)(\gamma)\) , i.e.~\(u_{\delta} = \operatorname{Int}(\delta)\) .

Let \(x, y, z\) be points of X. For \(\delta \in \varpi_{x,y}(X)\) and \(\eta \in \varpi_{y,z}(X)\) , we have \(u_{\delta \eta} = u_{\delta} \circ u_{\eta}\) , which is also written \(\delta \eta \gamma = {}^{\delta}(\eta \gamma)\) for \(\gamma \in \pi_1(X, z)\) .

Let \(x, y\) be points of X and let \(\delta\) , \(\delta'\) be elements of \(\varpi_{x,y}(\mathrm{X})\) . We have

\[
u _ {\delta^ {\prime}} = u _ {\delta} \circ \operatorname{Int} (\delta^ {- 1} \delta^ {\prime}) = \operatorname{Int} (\delta^ {\prime} \delta^ {- 1}) \circ u _ {\delta}.\tag{7}
\]

Remarks. --- 2) Let X be a topological space and let C be a path-connected component of X. If x and y are points of C, the topological space \(\Lambda_{x,y}(\mathrm{C})\) is identified with the topological space \(\Lambda_{x,y}(\mathrm{X})\) , so that the set \(\varpi_{x,y}(\mathrm{C})\) is identified with the set \(\varpi_{x,y}(\mathrm{X})\) . Thus, the fundamental group \(\varpi(\mathrm{C})\) is identified with the full subgroupoid of \(\varpi(\mathrm{X})\) having C as its set of points. In particular, for every point x of C, the group \(\pi_1(\mathrm{C}, x)\) is identified with the group \(\pi_1(\mathrm{X}, x)\) .

\begin{enumerate}
\def\labelenumi{\arabic{enumi})}
\setcounter{enumi}{2}
\tightlist
\item
  Let X be a topological space and let \(x\) , \(y\) , \(z\) be points of X. The map \((c,d)\mapsto c*d\) from \(\Lambda_{x,y}(\mathrm{X})\times \Lambda_{y,z}(\mathrm{X})\) to \(\Lambda_{x,z}(\mathrm{X})\) is continuous (III, p.~258, corollary of prop. 2). Note that the composition map \(\varpi_{x,y}(\mathrm{X})\times \varpi_{y,z}(\mathrm{X})\to \varpi_{x,z}(\mathrm{X})\) deduced from it is not necessarily continuous when one equips the sets \(\varpi_{x,y}(\mathrm{X})\) , \(\varpi_{y,z}(\mathrm{X})\) and \(\varpi_{x,z}(\mathrm{X})\) with quotient topologies (cf.~TG, I, p.~35 and III, p.~259). However, for every \(\gamma_0\in \varpi_{x,y}(\mathrm{X})\) and every \(\delta_0\in \varpi_{y,z}(\mathrm{X})\) , the partial mappings \(\gamma \mapsto \gamma \delta_0\) from \(\varpi_{x,y}(\mathrm{X})\) to \(\varpi_{x,z}(\mathrm{X})\) and \(\delta \mapsto \gamma_0\delta\) from \(\varpi_{y,z}(\mathrm{X})\) to \(\varpi_{x,z}(\mathrm{X})\) are homeomorphisms. Indeed, let \(c_0\in \Lambda_{x,y}(\mathrm{X})\) be a path of class \(\gamma_0\) . The map \(d\mapsto c_0*d\) is a continuous map from \(\Lambda_{y,z}(\mathrm{X})\) to \(\Lambda_{x,z}(\mathrm{X})\) . The map \(\delta \mapsto \gamma_0\delta\) follows by passing to quotients, hence is continuous. The same holds for the map \(\delta^{\prime}\mapsto \gamma_{0}^{-1}\delta^{\prime}\) from \(\varpi_{x,z}(\mathrm{X})\) to \(\varpi_{y,z}(\mathrm{X})\) . Since these two maps are inverses of each other, they are homeomorphisms. One argues analogously for the map \(\gamma \mapsto \gamma \delta_0\) . See also IV, p.~374.
\end{enumerate}

\subsection{Functoriality of the Poincaré groupoid}

Let X, Y be topological spaces and let \(f: X \to Y\) be a continuous map. The map \(c \mapsto f \circ c\) is a continuous map, denoted \(\Lambda(f)\) , from \(\Lambda(X) = \mathcal{C}_c(\mathbf{I}; X)\) to \(\Lambda(Y) = \mathcal{C}_c(\mathbf{I}; Y)\) (I, p.~132, lemma). It defines, by passing to subsets, the continuous maps

\[
\Lambda_ {x} (f) \colon \Lambda_ {x} (\mathrm{X}) \to \Lambda_ {f (x)} (\mathrm{Y}),
\]

continuous for \(x \in X\) ,

\[
\Lambda_ {x, y} (f) \colon \Lambda_ {x, y} (\mathrm{X}) \to \Lambda_ {f (x), f (y)} (\mathrm{Y}),
\]

continuous for \(x, y \in \mathrm{X}\) and

\[
\Omega (f) \colon \Omega (\mathrm{X}) \to \Omega (\mathrm{Y}).
\]

For \(x \in X\) , the map \(\Lambda_{x,x}(f)\) is also denoted \(\Omega_x(f)\) . By passing to arc-connected components (III, p.~290), we deduce from the map \(\Lambda_{x,y}(f)\) a map

\[
\varpi_ {x, y} (f) \colon \varpi_ {x, y} (\mathrm{X}) \to \varpi_ {f (x), f (y)} (\mathrm{Y}).
\]

Let \(x, y, z\) be points of X and let \(c \in \Lambda_{x,y}(X)\) and \(d \in \Lambda_{y,z}(X)\) be paths; by definition of the juxtaposition of paths, we have

\[
f \circ (c * d) = (f \circ c) * (f \circ d).
\]

By passing to strict homotopy classes, this relation follows.

\[
\varpi_ {x, z} (f) (\gamma \delta) = (\varpi_ {x, y} (f) (\gamma)) (\varpi_ {y, z} (f) (\delta))
\]

 for all \(\gamma \in \varpi_{x,y}(\mathrm{X})\) and every \(\delta \in \varpi_{y,z}(\mathrm{X})\) . Thus, the continuous map \(f\) and the maps \(\varpi_{x,y}(f)\) , for \(x\) and \(y \in \mathrm{X}\) , define a morphism of the groupoid \(\varpi(\mathrm{X})\) in the groupoid \(\varpi(\mathrm{Y})\) (II, p.~161, def. 3). It is denoted \(\varpi(f)\) and is called the Poincaré groupoid morphism deduced from the continuous map \(f\) . In particular, if \(x\) is a point of X, the map \(\varpi_{x,x}(f)\) is a homomorphism of the group \(\pi_1(\mathrm{X}, x)\) in the group \(\pi_1(\mathrm{Y}, f(x))\) ; this homomorphism is also denoted \(\pi_1(f, x)\) .

Remark 1. --- The homomorphism \(\varpi_{x,y}(f)\) is continuous if one equips the sets \(\varpi_{x,y}(\mathrm{X})\) and \(\varpi_{f(x),f(y)}(\mathrm{Y})\) with the quotient topology of the compact-open topology on \(\mathcal{C}_c(\mathbf{I};\mathrm{X})\) and \(\mathcal{C}_c(\mathbf{I};\mathrm{Y})\) .

To simplify the notation, if \(x\) , \(y \in X\) and \(\gamma \in \varpi_{x,y}(X)\) , we will sometimes write \(f_{*}(\gamma)\) for the element \(\varpi_{x,y}(f)(\gamma)\) of \(\varpi_{f(x),f(y)}(Y)\) .

Let X, Y, Z be topological spaces, let \(f: X \to Y\) and \(g: Y \to Z\) be continuous maps. For every path c in X, we have \((g \circ f) \circ c = g \circ (f \circ c)\) . It follows that one has

\[
\varpi (g \circ f) = \varpi (g) \circ \varpi (f).
\]

Let X and Y be topological spaces, let \(f_{0}\) and \(f_{1}\) be continuous maps from X to Y, and let \(\sigma\colon X\times I\to Y\) be a homotopy connecting \(f_{0}\) to \(f_{1}\) . The map \((c,t)\mapsto c(t)\) from \(\mathcal{C}_{c}(\mathbf{I};\mathrm{X})\times\mathbf{I}\) to X (III, p.~257, prop. 1) being continuous, the map from \(\mathcal{C}_{c}(\mathbf{I};\mathrm{X})\times\mathbf{I}\times\mathbf{I}\) to Y given by \((c,t,s)\mapsto\sigma(c(t),s)\) is also continuous. Consequently, the map \(\Sigma\colon(c,s)\mapsto\sigma(c(\cdot),s)\) is a continuous map from \(\mathcal{C}_{c}(\mathbf{I};\mathrm{X})\times\mathbf{I}\) to \(\mathcal{C}_{c}(\mathbf{I};\mathrm{Y})\) (loc. cit.). The map \(\Sigma\) is a homotopy connecting the maps \(\Lambda(f_{0})\) and \(\Lambda(f_{1})\) . By restriction to loop spaces, the map \(\Sigma\) induces a homotopy \(\Omega(\mathrm{X})\times\mathbf{I}\to\Omega(\mathrm{Y})\) connecting \(\Omega(f_{0})\) to \(\Omega(f_{1})\) . Let x be a point of X; suppose that \(\sigma\) is a pointed homotopy at x and set \(y=f_{0}(x)=f_{1}(x)\) . The map \(\Sigma\) then induces a continuous map from \(\Omega_{x}(\mathrm{X})\times\mathbf{I}\) to \(\Omega_{y}(\mathrm{Y})\) which is a pointed homotopy at \(e_{x}\) , connecting the maps \(\Omega_{x}(f_{0})\) and \(\Omega_{x}(f_{1})\) .

PROPOSITION 3. --- Let X and Y be topological spaces, \(f_{0}\) and \(f_{1}\) be continuous maps from X to Y, and let \(\sigma\colon\mathbf{X}\times\mathbf{I}\to\mathbf{Y}\) be a homotopy connecting \(f_{0}\) to \(f_{1}\) . Let x be a point of X; set \(y_{0}=f_{0}(x)\) , \(y_{1}=f_{1}(x)\) and denote by \(\delta\in\varpi_{y_{0},y_{1}}(\mathbf{Y})\) the class of the path d defined by \(d(t)=\sigma(x,t)\) for \(t\in\mathbf{I}\) . For all \(\gamma\in\pi_{1}(\mathbf{X},x)\) , we have \((f_{1})_{*}(\gamma)=\delta^{-1}\big((f_{0})_{*}(\gamma)\big)\delta\) .

Let c be a loop in X at x and let \(\gamma\) be its strict homotopy class. Set \(\gamma_{0} = (f_{0})_{*}(\gamma)\) and \(\gamma_{1} = (f_{1})_{*}(\gamma)\) ; these are the strict homotopy classes of \(f_{0} \circ c\) and \(f_{1} \circ c\) . For \((t, s) \in \mathbf{I} \times \mathbf{I}\) , set \(\varphi(t, s) = \sigma(c(t), s)\) . For all \(t \in I\) , we have \(\varphi(t, 0) = (f_{0} \circ c)(t)\) , \(\varphi(t, 1) = (f_{1} \circ c)(t)\) and \(\varphi(0, t) = \varphi(1, t) = d(t)\) . The relation \(\gamma_{0}\delta = \delta\gamma_{1}\) therefore follows from the following lemma.

Lemma 1. --- Let Y be a topological space and let \(\varphi\colon\mathbf{I}\times\mathbf{I}\to\mathbf{Y}\) be a continuous map. For \(t\in\mathbf{I}\) , set \(c_{0}(t)=\varphi(t,0)\) , \(c_{1}(t)=\varphi(t,1)\) , \(d_{0}(t)=\varphi(0,t)\) and \(d_{1}(t)=\varphi(1,t)\) . The paths \(c_{0}*d_{1}\) and \(d_{0}*c_{1}\) are strictly homotopic.

Let c denote the path in \(I \times I\) obtained by juxtaposing the paths \(t \mapsto (t,0)\) and \(t \mapsto (1,t)\) ; let d denote the path in \(I \times I\) obtained by juxtaposing the paths \(t \mapsto (0, t)\) and \(t \mapsto (t, 1)\) . The paths \(c\) and \(d\) have the same origin \((0, 0)\) and the same endpoint \((1, 1)\) . The map \((t, s) \mapsto (1 - s)c(t) + sd(t)\) from \(\mathbf{I} \times \mathbf{I}\) to \(\mathbf{I} \times \mathbf{I}\) is a strict homotopy joining \(c\) to \(d\) . We have \(c_0 * d_1 = \varphi \circ c\) and \(d_0 * c_1 = \varphi \circ d\) ; these two paths are therefore strictly homotopic.

COROLLARY 1. --- Let X and Y be topological spaces and let x be a point of X. Let \(f_0\) and \(f_1\) be continuous maps from X to Y. If there exists a pointed homotopy at x connecting \(f_0\) to \(f_1\) , we have \(\pi_1(f_0, x) = \pi_1(f_1, x)\) .

Let \(\sigma\) be a pointed homotopy at \(x\) connecting \(f_0\) to \(f_1\) . With the notations of Prop. 3, \(\delta\) is the class of a constant path with image \(f_0(x) = f_1(x)\) . The assertion follows from this.

COROLLARY 2. --- Let X and Y be topological spaces and let \(f\colon X \to Y\) be a homotopy equivalence. For every point \(x\) of X, the homomorphism

\[
\pi_ {1} (f, x) \colon \pi_ {1} (\mathrm{X}, x) \to \pi_ {1} (\mathrm{Y}, f (x))
\]

is an isomorphism.

Let \(g\) be a continuous map from Y to X, inverse up to homotopy of \(f\) . Let \(x\) be a point of X. It follows from prop. 3 applied to the homotopic maps \(\mathrm{Id}_{\mathrm{X}}\) and \(g \circ f: \mathrm{X} \to \mathrm{X}\) that the map \(\pi_1(g \circ f, x)\) is an isomorphism from the group \(\pi_1(\mathrm{X}, x)\) to the group \(\pi_1(\mathrm{X}, g \circ f(x))\) . Since

\[
\pi_ {1} (g \circ f, x) = \pi_ {1} (g, f (x)) \circ \pi_ {1} (f, x),
\]

the homomorphism \(\pi_{1}(f,x)\) is injective and the homomorphism \(\pi_{1}(g,f(x))\) is surjective. Since the map g is also a homotopy equivalence, the homomorphism \(\pi_{1}(g,f(x))\) is injective; it is therefore an isomorphism. Consequently, \(\pi_{1}(f,x)\) is an isomorphism.

Example. --- Let G be a topological group, and let e be its identity element. For every point g of G, the left and right translations, \(x \mapsto gx\) and \(x \mapsto xg\) are homeomorphisms of G onto itself (TG, III, p.~2) which map e to g. By Corollary 2, they induce isomorphisms of \(\pi_{1}(G, e)\) onto \(\pi_{1}(G, g)\) . These isomorphisms are not necessarily equal (IV, p.~459, exerc. 1).

COROLLARY 3. --- Let X be a topological space homotopy equivalent to a point. For every point x of X, the group \(\pi_{1}(X,x)\) is reduced to the identity element.

Cor. 3 applies in particular when X is the n-dimensional numerical space \(R^{n}\) and more generally when X is a subset of the space \(R^{n}\) which is star-shaped (III, p.~234) with respect to one of its points.

PROPOSITION 4. --- Let X be the product space of a family \((\mathrm{X}_{j})_{j\in\mathrm{J}}\) of topological spaces. The morphism of the groupoid \(\varpi(\mathrm{X})\) into the product of the groupoids \(\varpi(\mathrm{X}_{j})\) , for \(j\in J\) , defined by the family of morphisms \((\varpi(\mathrm{pr}_{j}))_{j\in\mathrm{J}}\) is an isomorphism.

Denote by \(\varphi\) this morphism of groupoids. The map deduced from it by passing to vertices is the identity map \(X \to \prod_{j} X_{j}\) . Let \(x = (x_{j})\) and \(y = (y_{j})\) be two points of \(X\) . Denote by \(\varphi_{x,y}\) the map from \(\varpi_{x,y}(X)\) to \(\prod_{j} \varpi_{x_{j},y_{j}}(X_{j})\) deduced from \(\varphi\) . If for every \(j \in J\) , \(c_{j}: I \to X_{j}\) is a path connecting \(x_{j}\) to \(y_{j}\) , the map \(t \mapsto (c_{j}(t))\) is a path in \(X\) connecting \(x\) to \(y\) (TG, I, p.~25, prop. 1). This proves that \(\varphi_{x,y}\) is surjective. Let \(c\) and \(d\) be two paths in \(X\) connecting \(x\) to \(y\) . Suppose that for every \(j \in J\) there exists a strict homotopy \(\sigma_{j}: I \times I \to X_{j}\) connecting \(\operatorname{pr}_{j} \circ c\) to \(\operatorname{pr}_{j} \circ d\) . Then, the map \((t,s) \mapsto (\sigma_{j}(t,s))\) from \(I \times I\) to \(X\) is a strict homotopy joining \(c\) to \(d\) (loc. cit.). This proves that \(\varphi_{x,y}\) is injective.

COROLLARY. --- Let \(x = (x_{j})_{j \in \mathrm{J}}\) be a point of X. The map

\[
\pi_ {1} (X, x) \to \prod_ {j \in J} \pi_ {1} (X _ {j}, x _ {j})
\]

deduced from the maps \(\pi_1(\mathrm{pr}_j, x_j)\) is a group isomorphism.

This isomorphism is said to be canonical. In what follows, we will often identify \(\pi_{1}(\mathrm{X},x)\) with \(\prod_{j\in\mathrm{J}}\pi_{1}(\mathrm{X}_{j},x_{j})\) by means of this isomorphism.

Remarks. --- 2) Let \((\mathrm{X}_{j})_{j\in\mathrm{J}}\) be a family of topological spaces. Let X denote the product topological space \(\prod_{j\in J}X_{j}\) and let \(x=(x_{j})\) be a point of X.

For every \(i \in \mathrm{J}\) , let \(u_i: \mathrm{X}_i \to \mathrm{X}\) be the map such that, for \(z \in \mathrm{X}_i\) , \(\operatorname{pr}_i \circ u_i(z) = z\) and, for all \(j \in \mathrm{J}\) distinct from \(i\) , \(\operatorname{pr}_j \circ u_i(z) = x_j\) . The map \(u_i\) is continuous and the map \(\pi_1(u_i, x_i)\) is identified with the canonical injection of the factor \(\pi_1(\mathrm{X}_i, x_i)\) in the product group of the family \((\pi_1(\mathrm{X}_j, x_j))_{j \in \mathrm{J}}\) (cf.~A, I, p.~45).

Suppose that the set J is finite and, for every \(j \in J\) , let \(\gamma_{j}\) be an element of \(\pi_{1}(X_{j}, x_{j})\) . The element \((\gamma_{j})\) of \(\pi_{1}(X, x)\) is the composite of the loop classes \((u_{j})_{*}(\gamma_{j})\) , \(j \in J\) , which are pairwise permutable.

\begin{enumerate}
\def\labelenumi{\arabic{enumi})}
\setcounter{enumi}{2}
\tightlist
\item
  Let \((\mathrm{X}_j)_{j\in \mathrm{J}}\) be a family of topological spaces. Let X denote the product topological space \(\prod_{j\in \mathrm{J}}\mathrm{X}_j\) and let \(x = (x_{j})\) be a point of X.
\end{enumerate}

Endow the sets \(\pi_1(\mathrm{X},x)\) and \(\pi_1(\mathrm{X}_j,x_j)\) with the quotient topology of the compact-open topology on the spaces \(\Lambda_x(\mathrm{X})\) and \(\Lambda_{x_j}(\mathrm{X}_j)\) . The isomorphism \(\pi_1(\mathrm{X},x)\to \prod_{j\in \mathrm{J}}\pi_1(\mathrm{X}_j,x_j)\) is then a homeomorphism. It is continuous (III, p.~294, remark 1). The compact-open topology on \(\Lambda (\mathrm{X})\) is generated by the subsets of the form \(T(\mathrm{K},\mathrm{U})\) , where K is a compact subset of I and U is an open subset of X. For \(j\in \mathrm{J}\) , let \(\mathrm{U}_j\) be an open subset of \(\mathrm{X}_j\) such that \(\prod_{j\in \mathrm{J}}\mathrm{U}_j\subset \mathrm{U}\) . Then \((\mathrm{pr}_j)_*(T(\mathrm{K},\mathrm{U}))\) contains \(T(\mathrm{K},\mathrm{U}_j)\) . This shows that the maps \((\mathrm{pr}_j)_*: \Lambda_x(\mathrm{X})\to \Lambda_{x_j}(\mathrm{X}_j)\) are open, and the maps \(\pi_1(\mathrm{pr}_j,x_j)\) are also open. Since they are surjective, the map \(\pi_1(\mathrm{X},x)\to \prod_{j\in \mathrm{J}}\pi_1(\mathrm{X}_j,x_j)\) is open (TG, I, p.~34, prop. 8). Being continuous and bijective, it is a homeomorphism (TG, I, p.~30, example 2).

PROPOSITION 5. — Let X be a topological space and let \((\mathrm{A}_{i})_{i\in\mathrm{I}}\) be an increasing family of subsets of X, indexed by a filtered ordered set I, such that every quasi-compact subset of X is contained in one of the \(A_{i}\) . The canonical groupoid morphism

\[
\rho \colon \varinjlim_ {i \in \mathrm{I}} \varpi (\mathrm{A} _ {i}) \to \varpi (\mathrm{X}),
\]

deduced from the canonical injections of \(A_{i}\) into X, is an isomorphism. If \(i \leqslant j\) , denote by \(\rho_{j,i}\) the groupoid morphism \(\varpi(\mathrm{A}_{i}) \to \varpi(\mathrm{A}_{j})\) deduced from the injection of \(A_{i}\) into \(A_{j}\) . Since the map deduced from \(\rho_{j,i}\) by passing to vertices is the injection \(A_{i} \to A_{j}\) and the \(A_{i}\) cover X, the map induced from \(\rho\) by passing to vertices is bijective.

Let \(a\) and \(b\) be points of X and let \(c\) be a path connecting \(a\) to \(b\) in X. The image of \(c\) is a quasi-compact subset of X (TG, I, p.~62, th. 2), since \(\mathbf{I}\) is compact. There therefore exists an element \(i \in \mathrm{I}\) such that the image of \(c\) is contained in \(\mathrm{A}_i\) . Consequently, the map deduced from \(\rho\) by passing to arrow sets is surjective.

Let \(i \in \mathrm{I}\) , let \(a\) and \(b\) be points of \(\mathrm{A}_i\) , let \(c\) and \(c'\) be paths joining \(a\) to \(b\) in \(\mathrm{A}_i\) ; let \(h\) be a strict homotopy joining \(c\) to \(c'\) in \(\mathrm{X}\) . Since \(\mathbf{I} \times \mathbf{I}\) is compact, \(h(\mathbf{I} \times \mathbf{I})\) is a quasi-compact subset of \(\mathrm{X}\) (loc. cit.) and there exists an element \(i \in \mathrm{I}\) such that the image of \(h\) is contained in \(\mathrm{A}_i\) . The paths \(c\) and \(c'\) are strictly homotopic in \(A_{i}\) ; a fortiori, the path classes {[}c{]} and \([c']\) have the same image in \(\varinjlim\varpi(A_{i})\) . Consequently, \(\rho\) is injective.

COROLLARY. --- Let a be a point of X and let J be the set of \(i \in I\) such that \(a \in A_{i}\) . The canonical homomorphism

\[
\varinjlim_ {i \in J} \pi_ {1} (A _ {i}, a) \to \pi_ {1} (X, a)
\]

is bijective.

Remark 4. --- The proposition and its corollary apply in particular when \((\mathrm{A}_{i})_{i\in\mathrm{I}}\) is an increasing family of subsets of X indexed by a filtered ordered set I such that the interiors of the \(A_{i}\) cover X.

\subsection{Freely homotopic loops}

DEFINITION 4. --- Let X be a topological space and let c and c' be two loops in X. One calls a free homotopy joining c to c' a homotopy σ joining c to c' such that σ(0, s) = σ(1, s) for all s ∈ I. One says that c is freely homotopic to c' if there exists a free homotopy joining c to c'.

The free homotopies joining \(c\) to \(c'\) correspond to the paths in the space \(\Omega(\mathrm{X})\) of loops of X joining \(c\) to \(c'\) . Consequently, the relation “ \(c\) is freely homotopic to \(c'\) ” is an equivalence relation on \(\Omega(\mathrm{X})\) whose equivalence classes are the arcwise connected components of \(\Omega(\mathrm{X})\) .

Note. --- Let \(\varphi\) be the canonical map from \(\mathbf{R}\) onto \(\mathbf{T} = \mathbf{R}/\mathbf{Z}\) (TG, V, p.~2). The map \(f \mapsto f \circ \varphi|\mathbf{I}\) is a homeomorphism of \(\mathcal{C}_c(\mathbf{T};\mathbf{X})\) onto \(\Omega(\mathbf{X})\) , whence, by passing to arcwise connected components, a bijection from the set {[}T; X{]} (III, p.~230) onto the set of free homotopy classes of loops in X.

PROPOSITION 6. --- Let X be an arcwise connected topological space and let x be a point of X.

\begin{enumerate}
\def\labelenumi{\alph{enumi})}
\item
  Every loop in X is freely homotopic to a loop at x. More precisely, if c is a loop at y and d is a path with origin y and endpoint x, c is freely homotopic to the loop \((\overline{d} * c) * d\) at x.
\item
  Two loops in X at x are freely homotopic if and only if their strict homotopy classes are conjugate in the group \(\pi_{1}(X,x)\) .
\end{enumerate}

Let us prove \(a\) ). For every \(s \in [0,1]\) , denote by \(d_s\) the path in X defined by \(d_s(t) = d(st)\) for \(t \in \mathbf{I}\) ; its origin is \(y\) . As the map \((s,t) \mapsto d(st)\) is continuous, the map \(s \mapsto d_s\) from \(\mathbf{I}\) to \(\mathcal{C}_c(\mathbf{I};\mathbf{X})\) is continuous (III, p.~257, prop. 1). The map \(s \mapsto (\overline{d_s} * c) * d_s\) is then a path in \(\Omega(\mathbf{X})\) (III, p.~257, prop. 2) joining \((e_y * c) * e_y\) to \((\overline{d} * c) * d\) , whence \(a\) ).

Let c and \(c'\) be two loops in X at x. If their strict homotopy classes are conjugate in \(\pi_{1}(X,x)\) , there exists a loop d at x such that \(c'\) is strictly homotopic to the loop \((\overline{d} * c) * d\) . It follows from a) that c and \(c'\) are freely homotopic. Conversely, suppose there exists a free homotopy \(\varphi\) joining c to \(c'\) . Set \(d(t) = \varphi(0,t)\) , we also have \(d(t) = \varphi(1,t)\) and d is a loop at x. By lemma 1 of III, p.~295, the loops \(c * d\) and \(d * c'\) are strictly homotopic. The strict homotopy classes of c and \(c'\) are therefore conjugate in \(\pi_{1}(X,x)\) .

Scholium. --- Let X be an arcwise connected topological space and let x be a point of X. Proposition 6 allows one to define a canonical bijection from the set of free homotopy classes of loops in X onto the set of conjugacy classes in \(\pi_{1}(X,x)\) .

\section{HOMOTOPY AND COVERINGS}

\subsection{Homotopy and coverings}

PROPOSITION 1. --- Let B be a topological space and let E be a covering of B. Let B' be a topological space and let \(f_0\) and \(f_1\) be continuous maps from B' to B. If the maps \(f_0\) and \(f_1\) are homotopic, the coverings \(f_0^*(\mathrm{E})\) and \(f_1^*(\mathrm{E})\) of B' are isomorphic.

Let \(\sigma\colon\mathrm{B}^{\prime}\times\mathbf{I}\to\mathrm{B}\) be a homotopy connecting \(f_{0}\) to \(f_{1}\) . Denote by \(i_{0}\) and \(i_{1}\) the maps \(x\mapsto(x,0)\) and \(x\mapsto(x,1)\) from \(\mathrm{B}^{\prime}\) to \(\mathrm{B}^{\prime}\times\mathbf{I}\) . If \(t\in\{0,1\}\) , we have \(f_{t}=\sigma\circ i_{t}\) and the coverings \(f_{t}^{*}(\mathrm{E})\) and \(i_{t}^{*}(\sigma^{*}(\mathrm{E}))\) of \(\mathrm{B}^{\prime}\) are therefore isomorphic (I, p.~15). Since the topological space \(\mathbf{I}\) is locally connected and simply connected (I, p.~127, corollary), the coverings \(i_{0}^{*}(\sigma^{*}(\mathrm{E}))\) and \(i_{1}^{*}(\sigma^{*}(\mathrm{E}))\) of \(B'\) are isomorphic (I, p.~130, cor. 1 of prop. 8).

COROLLARY. --- A topological space homeomorphic to a simply connected topological space is simply connected.

Let B and B' be topological spaces, let \(f: B \rightarrow B'\) be a homotopy equivalence and let \(g: B' \rightarrow B\) be a continuous map, inverse up to homotopy of f. Let E be a covering of B. Since \(g \circ f\) is homotopic to the identity map of B, the covering E is isomorphic to the covering \(f^{*}(g^{*}(\mathrm{E}))\) (prop. 1). If the space B' is simply connected, the covering \(g^{*}(\mathrm{E})\) is trivializable. Consequently, the covering E is trivializable. This proves that the space B is simply connected.

Remark. --- Let G be a discrete topological group. With the notation of the proposition, suppose that E is a principal covering with group G. The coverings \(f_{0}^{*}(E)\) and \(f_{1}^{*}(E)\) are then isomorphic principal coverings (I, p.~131, remark).

PROPOSITION 2 (Lifting of homotopies). --- Let B be a topological space and let E be a covering of B. Let B' be a topological space and let \(f_0\) and \(f_1\) be continuous maps from B' to B. Let \(\sigma: \mathrm{B}' \times \mathbf{I} \to \mathrm{B}\) be a homotopy connecting \(f_0\) to \(f_1\) . Let \(g_0: \mathrm{B}' \to \mathrm{E}\) be a continuous lifting of \(f_0\) to E. There exists a unique continuous map \(\widetilde{\sigma}: \mathrm{B}' \times \mathbf{I} \to \mathrm{E}\) lifting \(\sigma\) which is a homotopy with origin \(g_0\) . Its endpoint \(g_1: \mathrm{B}' \to \mathrm{E}\) is a continuous lifting of \(f_1\) to E.

This follows from corollary 3 of prop. 8 of I, p.~130, applied to the simply connected and locally connected topological space I.

\subsection{Lifting of paths}

PROPOSITION 3. --- Let B be a topological space and let \(p\colon E \to B\) be a covering. Denote by \(\widetilde{p}\colon \mathcal{C}_{c}(\mathbf{I};E) \to \mathcal{C}_{c}(\mathbf{I};B)\) the map \(c \mapsto p \circ c\) . Denote by \(o_{\mathrm{E}}\colon \mathcal{C}_{c}(\mathbf{I};E) \to E\) (resp. \(o_{\mathrm{B}}\colon \mathcal{C}_{c}(\mathbf{I};B) \to B\) ) the map defined by \(c \mapsto c(0)\) . Then the diagram

\[
\begin{array}{c} \mathcal {C} _ {c} (\mathbf {I}; \mathrm{E}) \xrightarrow {o _ {\mathrm{E}}} \mathrm{E} \\ \Big \downarrow \widetilde {p} \\ \mathcal {C} _ {c} (\mathbf {I}; \mathrm{B}) \xrightarrow {o _ {\mathrm{B}}} \mathrm{B} \end{array}
\]

is a Cartesian square.

The topological space I is locally connected, locally compact, and simply connected. The proposition thus follows from prop. 9 of I, p.~131, applied with T = I and t = 0.

COROLLARY 1. --- The map \(\widetilde{p}\colon\mathcal{C}_{c}(\mathbf{I};\mathrm{E})\to\mathcal{C}_{c}(\mathbf{I};\mathrm{B})\) is a covering.

This follows from I, p.~71, cor. 2 of prop. 1.

COROLLARY 2 (Lifting of paths). --- Let \(p\) : E → B be a covering, let x be a point of E and a = p(x). The map c ↦ p ∘ c is a homeomorphism from the space \(\Lambda_x\) (E) of paths in E with origin x onto the space \(\Lambda_a\) (B) of paths in B with origin a.

With the notations of proposition 3, we have \(\Lambda_{a}(\mathrm{B}) = o_{\mathrm{B}}^{-1}(a)\) and \(\Lambda_x(\mathrm{E}) = o_{\mathrm{E}}^{-1}(x)\) . The corollary then follows from I, p.~10, cor. of prop. 4.

We shall say that a continuous map \(p \colon \mathrm{E} \to \mathrm{B}\) satisfies the path-lifting property if, for every path \(c \colon \mathbf{I} \to \mathbf{B}\) and every point \(x\) of E lifting \(c(0)\) , there exists a path \(c'\) in E with origin \(x\) which lifts \(c\) . Such a path is then unique if \(p\) is étale and separated (I, p.~34, cor. 1 of prop. 11). Let E be a covering and \(p\) its projection; the map \(p\) is étale, separated, and possesses the path-lifting property (corollary 2). For a partial converse, cf.~III, p.~315, corollary.

PROPOSITION 4. --- Let \(p \colon \mathrm{E} \to \mathrm{B}\) be an étale and separated map which satisfies the path-lifting property. Let \(a\) and \(b\) be points of \(\mathrm{B}\) and let \(x\) be a point of \(\mathrm{E}\) such that \(p(x) = a\) . Let \(c_0\) and \(c_1\) be two strictly homotopic paths joining \(a\) to \(b\) in \(\mathrm{B}\) . The paths with origin \(x\) in \(\mathrm{E}\) which respectively lift \(c_0\) and \(c_1\) have the same endpoint and are strictly homotopic.

Let \(\sigma\colon\mathbf{I}\times\mathbf{I}\to\mathbf{B}\) be a strict homotopy joining \(c_{0}\) to \(c_{1}\) . For \(s\in\mathbf{I}\) , let \(c_{s}^{\prime}\) be the unique path with origin \(x\) in E which lifts the path \(t\mapsto\sigma(t,s)\) . For \((t,s)\in\mathbf{I}\times\mathbf{I}\) , set \(\sigma^{\prime}(t,s)=c_{s}^{\prime}(t)\) ; we have \(p\circ\sigma^{\prime}=\sigma\) . The map \(\sigma^{\prime}\) is constant on \(\{0\}\times\mathbf{I}\) ; by construction, it is continuous on \(\mathbf{I}\times\{s\}\) for all \(s\in\mathbf{I}\) . It is therefore continuous, by I, p.~37, cor. 1 of th. 1. In particular, the map from \(\mathbf{I}\) to \(\mathrm{E}_{b}\) given by \(s\mapsto\sigma^{\prime}(1,s)\) is a continuous lifting of the constant path with image \(b\) ; it is therefore constant. This proves that the map \(\sigma^{\prime}\) is a strict homotopy joining \(c_{0}^{\prime}\) to \(c_{1}^{\prime}\) .

COROLLARY 1. --- Let B be a topological space and let E be a covering of B; denote by p its projection. For every pair \((x,y)\) of points of E, the map \(\varpi_{x,y}(p)\colon\varpi_{x,y}(\mathrm{E})\to\varpi_{p(x),p(y)}(\mathrm{B})\) is injective. In particular, for every point x of E, the homomorphism \(\pi_{1}(p,x)\colon\pi_{1}(\mathrm{E},x)\to\pi_{1}(\mathrm{B},p(x))\) is injective.

COROLLARY 2. --- Let B be a topological space and let E be a covering of B; denote by p its projection. Let x be a point of E, set \(b = p(x)\) . For the class in \(\pi_{1}(B, b)\) of a loop c of B at a to belong to the image subgroup of the homomorphism \(\pi_{1}(p, x)\) , it is necessary and sufficient that the path c' with origin x lifting c be a loop of E at x.

The condition is obviously sufficient. Conversely, let \(c_{1}^{\prime}\) be a loop of E at x. Set \(c_{1} = p \circ c_{1}^{\prime}\) and suppose that the loops c and \(c_{1}\) are strictly homotopic. By prop. 4, the path \(c^{\prime}\) has the same endpoint as the path \(c_{1}^{\prime}\) . It is therefore a loop at x.

COROLLARY 3. --- Let B be a topological space and let \((\mathrm{E}_{1}, p_{1})\) and \((\mathrm{E}_{2}, p_{2})\) be coverings of B. Denote \((\mathrm{E}, p)\) their B-product space and let \(x = (x_{1}, x_{2})\) be a point of E. The image of the homomorphism \(\pi_{1}(p, x)\) is the intersection of the images of the homomorphisms \(\pi_{1}(p_{1}, x_{1})\) and \(\pi_{1}(p_{2}, x_{2})\) .

Let \(c\) be a loop in B based at \(p_{1}(x_{1}) = p_{2}(x_{2})\) and, for \(i \in \{1, 2\}\) , let \(c_{i}^{\prime}\) be the path with origin \(x_i\) in \(E_i\) lifting \(c\). The B-space \(E = E_{1} \times_{B} E_{2}\) is a covering of B (I, p.~72, cor. 3 of prop. 2), and the path \(c^{\prime}: t \mapsto (c_{1}^{\prime}(t), c_{2}^{\prime}(t))\) is the unique path with origin \(x\) in E lifting \(c\). For \(c^{\prime}\) to be a loop, it is necessary and sufficient that \(c_{1}^{\prime}\) and \(c_{2}^{\prime}\) be loops. Corollary 3 thus follows from Corollary 2.

\subsection{Operations of the Poincaré groupoid in covering spaces}

Let B be a topological space and let \(p \colon \mathrm{E} \to \mathrm{B}\) be a separated étale map satisfying the path-lifting property (III, p.~302); this is the case if the map \(p\) makes E a covering of B (loc. cit.). Let \(b\) and \(b'\) be points of B and let \(c \in \Lambda_{b,b'}(\mathrm{B})\) be a path in B connecting \(b\) to \(b'\) . For every point \(x\) of the fiber \(\mathrm{E}_b\) , denote by \(x \cdot c\) the endpoint of the path with origin \(x\) in E that lifts \(c\) . The point \(x \cdot c\) belongs to the fiber \(E_{b'}\) and depends only on the class \(\gamma \in \varpi_{b,b'}(\mathrm{B})\) of the path c (III, p.~302, prop. 4); one will thus write \(x \cdot \gamma\) instead of \(x \cdot c\) .

If \(c\) is the constant path at \(b\) , we have \(x \cdot \gamma = x\) , because the constant path with origin \(x\) lifts \(c\) .

Let \(b''\) be another point of B and let \(c' \in \Lambda_{b',b''}(B)\) . For every point \(x\) of \(E_b\) , we have:

\[
(x \cdot c) \cdot c ^ {\prime} = x \cdot (c * c ^ {\prime}).\tag{1}
\]

Indeed, denote by \(\widetilde{c}\) the lifting with origin \(x\) of \(c\) and \(\widetilde{c}'\) the lifting of \(c'\) with origin \(x \cdot c = \widetilde{c}(1)\) . The paths \(\widetilde{c}\) and \(\widetilde{c}'\) are juxtaposable and \(\widetilde{c} * \widetilde{c}'\) is the path with origin \(x\) lifting \(c * c'\) .

Denote by \(\varphi_{b,b'}\colon \varpi_{b,b'}(\mathrm{B})\to \mathcal{F}(\mathrm{E}_b;\mathrm{E}_{b'})\) the map such that, for \(\gamma \in \varpi_{b,b'}(\mathrm{B})\) and \(x\in \mathrm{E}_b\) , we have

\[
\varphi_ {b, b ^ {\prime}} (\gamma) (x) = x \cdot \gamma .
\]

For every \(\gamma \in \varpi_{b,b'}(B)\) and every \(\gamma' \in \varpi_{b',b''}(B)\) , it follows from relation (1) that one has

\[
\varphi_ {b, b ^ {\prime \prime}} (\gamma \gamma^ {\prime}) = \varphi_ {b ^ {\prime}, b ^ {\prime \prime}} (\gamma^ {\prime}) \circ \varphi_ {b, b ^ {\prime}} (\gamma).\tag{2}
\]

The family \(\varphi = (\varphi_{b,b'})_{(b,b')\in\mathrm{B}\times\mathrm{B}}\) therefore defines a right operation law of the groupoid \(\varpi(\mathrm{B})\) on the set E with respect to the mapping \(p\colon\mathrm{E}\to\mathrm{B}\) (II, p.~167). It is called the canonical operation of the groupoid \(\varpi(\mathrm{B})\) associated with the mapping \(p\colon\mathrm{E}\to\mathrm{B}\) . The map \(\varphi_{b,b}\colon\pi_{1}(\mathrm{B},b)\to\mathcal{F}(\mathrm{E}_{b};\mathrm{E}_{b})\) defines a right action of the group \(\pi_{1}(\mathrm{B},b)\) on the fiber \(\mathrm{E}_{b}\) of E. This operation is called the canonical operation of \(\pi_{1}(\mathrm{B},b)\) on the fiber \(\mathrm{E}_{b}\) .

Remark. --- Let \(b\) and \(b'\) be two points of B, \(x\) be a point of \(\mathrm{E}_b\) , \(c \in \Lambda_{b,b'}(\mathrm{B})\) be a path in B and \(\widetilde{c}\) the path in E with origin \(x\) which lifts \(c\) . For all \(s \in \mathbf{I}\) , denote by \(c_s\) the path \(t \mapsto c(st)\) . The path in E with origin \(x\) lifting \(c_s\) is the path \(t \mapsto \widetilde{c}(st)\) , whose endpoint is the point \(\widetilde{c}(s)\) . We therefore have \(\widetilde{c}(s) = x \cdot c_s\) for all \(s \in \mathbf{I}\) .

Let B' be a topological space and let \(p'\colon E'\to B'\) be an étale and separated map which satisfies the path-lifting property. Let \(f\colon B\to B'\) and \(g\colon E\to E'\) be continuous maps such that \(p'\circ g=f\circ p\) . For \(b, b'\in B, \gamma\in\varpi_{b,b'}(B)\) and \(x\in E_b\) , we have:

\[
g (x \cdot \gamma) = g (x) \cdot f _ {*} (\gamma).\tag{3}
\]

Indeed, let c be a path in B whose strict homotopy class is \(\gamma\) ; denote by \(\widetilde{c}\) the path in E with origin x that lifts c. The path \(g \circ \widetilde{c}\) is then the lifting with origin \(g(x)\) of \(f \circ c\) in E'.

In particular, for \(B = B'\) and \(f = \operatorname{Id}_{B}\) , we have

\[
g (x \cdot \gamma) = g (x) \cdot \gamma .\tag{4}
\]

Let \(p \colon \mathrm{E} \to \mathrm{B}\) and \(p' \colon \mathrm{E}' \to \mathrm{B}\) be étale and separated maps satisfying the path-lifting property. Let \(b\) be a point of B. If \(f \colon \mathrm{E} \to \mathrm{E}'\) is a B-morphism, the map \(f_b \colon \mathrm{E}_b \to \mathrm{E}_b'\) is a \(\pi_1(\mathrm{B}, b)\)-morphism.

THEOREM 1. --- Let B be a topological space, let (E,p) be a covering of B, let b be a point of B and let x be a point of the fiber \(E_{b}\) .

\begin{enumerate}
\def\labelenumi{\alph{enumi})}
\item
  The orbit of x under the canonical action of the group \(\pi_{1}(\mathrm{B},b)\) is the intersection of \(E_{b}\) with the path-connected component of x in E. In particular, if the space E is path-connected, this operation is transitive.
\item
  The fixator (A, I, p.~51) of x is the subgroup \(p_{*}(\pi_{1}(E,x))\) of \(\pi_{1}(B,b)\) .
\item
  For every \(\gamma \in \pi_1(\mathrm{B},b)\) , we have \(p_*(\pi_1(\mathrm{E},x)) = \operatorname{Int}(\gamma)(p_*(\pi_1(\mathrm{E},x \cdot \gamma)))\) .
\item
  If the space B is connected and the covering E is Galois, the subgroup \(p_{*}(\pi_{1}(\mathrm{E},x))\) is distinguished in \(\pi_{1}(\mathrm{B},b)\) .
\end{enumerate}

Assertion a) is immediate. Assertion b) follows from cor. 2 of prop. 4 (III, p.~303). Assertion c) follows from b) and from proposition 2 of A, I, p.~52. Finally, suppose that B is connected and that E is a Galois covering of B. For every \(\gamma \in \pi_1(B,b)\) , there exists a B-automorphism \(g\) of E such that \(g(x) = x \cdot \gamma\) (I, p.~102, th. 2, c)). We have \(p = p \circ g\) , whence \(p_*(\pi_1(E,x)) = p_*(\pi_1(E,x \cdot \gamma))\) . Assertion d) therefore follows from c).
